% ERS_SRS_2B_v2.0.tex
% Especificacion de Requisitos de Software, AgroMoreira (SGA).
% Estandar base: ISO/IEC/IEEE 29148:2018 y ISO/IEC 25010:2023.
% Enfoque metodologico: legal-first (Amaral et al., 2021).
% Compilar: pdflatex -> bibtex -> pdflatex -> pdflatex

\documentclass[11pt,a4paper]{article}

\usepackage[utf8]{inputenc}
\usepackage[T1]{fontenc}
\usepackage[spanish,es-noshorthands]{babel}
\usepackage[a4paper,margin=2.4cm]{geometry}
\usepackage{longtable}
\usepackage{array}
\usepackage{booktabs}
\usepackage{enumitem}
\usepackage{fancyhdr}
\usepackage{titlesec}
\usepackage[hidelinks]{hyperref}
\usepackage{parskip}

\newcolumntype{L}[1]{>{\raggedright\arraybackslash}p{#1}}

\setlist{nosep,leftmargin=1.4em}
\titleformat{\section}{\large\bfseries}{\thesection}{0.6em}{}
\titleformat{\subsection}{\normalsize\bfseries}{\thesubsection}{0.5em}{}

\newcommand{\ruta}[1]{\texttt{#1}}
\newcommand{\rf}[2]{\subsection*{#1. #2}\addcontentsline{toc}{subsection}{#1. #2}}

\pagestyle{fancy}
\fancyhf{}
\fancyhead[L]{\small ERS/SRS AgroMoreira}
\fancyhead[R]{\small Version 2.0}
\fancyfoot[C]{\small Pagina \thepage}
\renewcommand{\headrulewidth}{0.4pt}

\begin{document}

\begin{titlepage}
\centering
\vspace*{2cm}
{\LARGE\bfseries Especificacion de Requisitos de Software\par}
\vspace{0.4cm}
{\large (ERS / SRS)\par}
\vspace{1.2cm}
{\Large\bfseries AgroMoreira\par}
{\large Sistema de Gestion Agricola con Inteligencia Artificial\par}
\vspace{2cm}
\begin{tabular}{rl}
\textbf{Version:} & 2.0 \\[2pt]
\textbf{Estandar base:} & ISO/IEC/IEEE 29148:2018 y ISO/IEC 25010:2023 \\[2pt]
\textbf{PFC:} & \#11, Categoria C, riesgo minimo operativo \\[2pt]
\textbf{Paralelo:} & 4to Software A \\[2pt]
\textbf{Fecha:} & 1 de septiembre de 2026 \\
\end{tabular}
\vspace{1.6cm}

\begin{tabular}{rl}
\textbf{Equipo:} & Jeanpierre Robinson Espinoza, lider \\
 & Danela Dayana Arteaga Alava, secretaria \\
 & Kamila Annabella Calle Delgado, documentadora \\
 & Maria del Rosario Escudero Plaza, modeladora \\
 & Roselyn Andreina Sanchez Centeno, verificadora \\[6pt]
\textbf{Docente supervisor:} & Ing. Gleiston Ciceron Guerrero Ulloa, PhD \\
\end{tabular}
\vfill
{\small Universidad Tecnica Estatal de Quevedo \textperiodcentered{} Facultad de Ciencias de la Computacion \textperiodcentered{} Ingenieria de Requerimientos (ISR-401)}
\end{titlepage}

\tableofcontents
\newpage

\section*{Historial de versiones}
\addcontentsline{toc}{section}{Historial de versiones}

\begin{longtable}{L{1.2cm} L{2.4cm} L{9.5cm} L{2.6cm}}
\toprule
\textbf{Version} & \textbf{Fecha} & \textbf{Cambios} & \textbf{Autor} \\
\midrule
\endfirsthead
\toprule
\textbf{Version} & \textbf{Fecha} & \textbf{Cambios} & \textbf{Autor} \\
\midrule
\endhead
1.0 & 30 de agosto de 2026 & Esqueleto inicial del documento, sin trabajo de campo todavia. & Equipo AgroMoreira \\
\addlinespace
1.1 & 31 de agosto de 2026 & Correccion del equipo y del enfoque metodologico a partir del expediente etico real. & Equipo AgroMoreira \\
\addlinespace
1.2 & 1 de septiembre de 2026 & Correccion del liderazgo del equipo. & Equipo AgroMoreira \\
\addlinespace
1.3 & 1 de septiembre de 2026 & Definicion del enfoque legal-first como enfoque oficial del proyecto. & Equipo AgroMoreira \\
\addlinespace
2.0 & 1 de septiembre de 2026 & Cierre de la Seccion 3 con los 39 requisitos funcionales y los 21 no funcionales obtenidos en las 17 entrevistas y las 7 sesiones de validación, y de las historias de usuario con criterios de aceptacion para los 17 requisitos Must have. Cierre del modelado del sistema en la Seccion 4. Matriz de trazabilidad cerrada en 60 filas y Seccion 6 con el producto minimo viable funcional. & Equipo AgroMoreira \\
\bottomrule
\end{longtable}

\newpage

\section{Introduccion}

\subsection{Proposito}

Este documento especifica los requisitos funcionales y no funcionales del Sistema de Gestion Agricola con Inteligencia Artificial para Agricola Moreira, finca de policultivo dedicada al cacao y al platano verde. El componente empirico del proyecto sigue el enfoque legal-first de Amaral, Abualhaija, Sabetzadeh y Briand \cite{amaral2021}. A partir de un modelo conceptual de cumplimiento legal organizado en tres bloques normativos, la Ley Organica de Proteccion de Datos Personales del Ecuador, la Resolucion 183 de AGROCALIDAD sobre buenas practicas y trazabilidad del cacao, y la Resolucion 0072 sobre bioseguridad y manejo fitosanitario, se definieron 26 criterios de cumplimiento documentados en \ruta{Modelo\_Legal\_LOPDP.md}. El trabajo de campo permitio evaluar cuales de esos criterios quedaban cubiertos por los requisitos obtenidos mediante entrevistas convencionales y cuales solo se cubren cuando se derivan requisitos directamente del texto legal.

\subsection{Alcance}

El sistema cubre la centralizacion del registro de parcelas, cultivos y variedades, el registro de actividades agricolas y cosechas, la gestion del inventario de insumos, la asignacion y el seguimiento de tareas, la generacion de reportes de produccion y costos, y un modulo de alertas de plagas asistido por inteligencia artificial que siempre requiere confirmacion humana antes de aplicarse a un lote. Tambien cubre los requisitos de cumplimiento legal derivados del modelo de proteccion de datos y de la normativa agroindustrial ecuatoriana. El alcance definitivo se fijo tras las 17 entrevistas y las 7 sesiones de validacion descritas en la Seccion 3, y ya fue confirmado con el administrador Raul Moreira Alay.

\subsection{Glosario de terminos}

\begin{longtable}{L{3.2cm} L{12cm}}
\toprule
\textbf{Termino} & \textbf{Definicion} \\
\midrule
\endhead
RF & Requisito funcional. \\
\addlinespace
RNF & Requisito no funcional. \\
\addlinespace
RD & Requisito derivado directamente del texto legal, sin evidencia de entrevista. \\
\addlinespace
Sugerencia explicable & Recomendacion generada por el componente de inteligencia artificial, acompanada de una justificacion comprensible para el usuario y de los datos en los que se basa, sin que el sistema ejecute la accion de forma autonoma. \\
\addlinespace
AGROCALIDAD & Agencia de Regulacion y Control Fito y Zoosanitario del Ecuador. \\
\addlinespace
LOPDP & Ley Organica de Proteccion de Datos Personales del Ecuador. \\
\addlinespace
Perfil tecnico & Usuario con familiaridad previa en el manejo de sistemas informaticos, como el personal de administracion o los tecnicos agronomos. \\
\addlinespace
Perfil no tecnico & Usuario sin familiaridad previa en sistemas informaticos, como buena parte del personal de campo. \\
\bottomrule
\end{longtable}

\subsection{Referencias}

Las fuentes citadas en este documento estan en \ruta{referencias.bib}, con verificacion cruzada de autor, titulo, publicacion, ano y DOI. Esta bibliografia no sustituye a la del protocolo de investigacion.

\subsection{Vision general del documento}

La Seccion 2 describe el producto y sus interesados. La Seccion 3 presenta los 39 requisitos funcionales, los 21 requisitos no funcionales y las historias de usuario con sus criterios de aceptacion, todos obtenidos del trabajo de campo real. La Seccion 4 recoge el modelado del sistema, con los casos de uso, los diagramas de clases, de comportamiento y de componentes, y los mockups. Las Secciones 5 y 6 describen la priorizacion, la trazabilidad cerrada y el producto minimo viable.

\section{Descripcion general}

\subsection{Perspectiva del producto}

AgroMoreira es un sistema nuevo e independiente, desarrollado especificamente para Agricola Moreira. No reemplaza un sistema legado formal, ya que la finca lleva sus registros actuales en papel y en hojas de calculo, segun se confirmo durante las entrevistas.

\subsection{Diagrama de contexto}

\begin{verbatim}
                +---------------------------------------+
 Propietario /  |                                       |   Tecnicos de
 Administrador -+                                       +-- AGROCALIDAD
 de la finca    |                                       |   (externo)
                |                                       |
 Ingeniero      |         SISTEMA AGROMOREIRA           |   Centro de acopio /
 agronomo    ---+   (gestion agricola + sugerencias     +-- Asociacion de
 residente      |    de IA explicables)                 |   productores (externo)
                |                                       |
 Capataces   ---+                                       |
 Personal de    |                                       |
 campo       ---+                                       |
 Personal de    |                                       |
 acopio      ---+                                       |
                +---------------------------------------+
\end{verbatim}

La version trazada de este diagrama es \ruta{03\_Modelado/Diagrams/CTX01\_Context\_Diagram.drawio}, y la matriz de poder e interes es \ruta{CTX02\_Power\_Interest\_Matrix.drawio}.

\subsection{Mapa de interesados, matriz de poder e interes}

\begin{longtable}{L{3.4cm} L{1.4cm} L{1.4cm} L{2.6cm} L{5.4cm}}
\toprule
\textbf{Interesado} & \textbf{Poder} & \textbf{Interes} & \textbf{Cuadrante} & \textbf{Rol en el proyecto} \\
\midrule
\endhead
Propietario o administrador de la finca & Alto & Alto & Gestionar de cerca & Decide la adopcion del sistema y firmo el aval institucional. \\
\addlinespace
Ingeniero agronomo residente & Medio & Alto & Mantener informado e involucrar & Informante clave tecnico y usuario del modulo de IA. \\
\addlinespace
Capataces & Medio & Medio & Mantener satisfecho & Coordina al personal de campo y usa el sistema como usuario intermedio. \\
\addlinespace
Personal de campo & Bajo & Alto & Mantener informado & Usuario final principal de las sugerencias explicables. \\
\addlinespace
Personal de acopio & Bajo & Medio & Monitorear & Usuario del modulo de trazabilidad y cosecha. \\
\addlinespace
Tecnicos de AGROCALIDAD, externo & Medio & Bajo & Monitorear & Informante externo sobre requisitos normativos. \\
\addlinespace
Centro de acopio o asociacion, externo & Bajo & Bajo & Monitorear & Informante externo sobre trazabilidad posterior a la cosecha. \\
\bottomrule
\end{longtable}

\subsection{Modelado organizacional i estrella}

El propietario o administrador depende del personal de campo para la ejecucion correcta de las labores culturales. El personal de campo depende del sistema AgroMoreira para recibir sugerencias comprensibles y accionables sobre el manejo del cultivo. El propio sistema depende del ingeniero agronomo para validar la calidad de esas sugerencias de inteligencia artificial. El propietario o administrador depende tambien del sistema para la trazabilidad y los reportes que debe presentar ante AGROCALIDAD. El personal de acopio depende del sistema para registrar la cosecha y la calidad posterior a esta.

Los diagramas de dependencia estrategica y de razon estrategica estan en \ruta{03\_Modelado/Diagrams/ISTAR01\_Strategic\_Dependency\_SD.drawio} y \ruta{ISTAR02\_Strategic\_Rationale\_SR.drawio}.

\subsection{Entorno operativo}

El personal de campo trabaja principalmente con telefonos Android de gama baja y con poca experiencia previa en aplicaciones moviles, segun se confirmo en las entrevistas a los jornaleros. La conectividad dentro de las parcelas es intermitente, lo que llevo a definir el requisito no funcional de robustez de campo de la Seccion 3.3. El personal tecnico y administrativo cuenta con acceso a computadores para las tareas de configuracion, reporte y analisis.

\subsection{Restricciones de diseno}

El sistema debe operar con conectividad intermitente o nula en el campo, tal como confirmo el personal de campo entrevistado. Las sugerencias de inteligencia artificial deben ser explicables y en ningun caso pueden ejecutar acciones de forma autonoma sobre un lote sin confirmacion humana. El sistema debe cumplir con la Ley Organica de Proteccion de Datos Personales del Ecuador \cite{lopdp2021}, segun se detalla a continuacion.

\subsection{Mapeo a la Ley Organica de Proteccion de Datos Personales del Ecuador}

Fuente legal verificada, Registro Oficial, Quinto Suplemento N.\textsuperscript{o} 459, del 26 de mayo de 2021.

\begin{longtable}{L{2.4cm} L{4cm} L{9cm}}
\toprule
\textbf{Articulo} & \textbf{Tema} & \textbf{Relevancia para AgroMoreira} \\
\midrule
\endhead
Art. 7 y 8 & Tratamiento legitimo y consentimiento & El consentimiento debe ser libre, especifico, informado e inequivoco, y revocable en cualquier momento. Es la base del proceso de entrevistas y del formulario de consentimiento informado. \\
\addlinespace
Art. 10, literal e & Pertinencia y minimizacion de datos & El sistema y el estudio solo recolectan los datos estrictamente necesarios para la gestion agricola. \\
\addlinespace
Art. 10, literal g & Confidencialidad & El tratamiento no puede comunicarse para un fin distinto al recogido. Respalda el compromiso de confidencialidad firmado por el equipo. \\
\addlinespace
Art. 10, literal j & Seguridad de datos personales & Respalda el cifrado AES-256 de la zona restringida de evidencias. \\
\addlinespace
Art. 13 a 19 & Derechos de acceso, rectificacion, eliminacion, oposicion y portabilidad & El sistema en su version productiva debe soportar estos cinco derechos, recogidos en el RF-23. \\
\addlinespace
Art. 20 & Derecho a no ser objeto de una decision basada unica o parcialmente en valoraciones automatizadas & Es el articulo mas relevante para AgroMoreira. El titular puede exigir una explicacion motivada sobre cualquier sugerencia del sistema, conocer los datos usados e impugnarla. Da base legal directa a los requisitos de explicabilidad del RF-09 y de la Seccion 3.3. \\
\bottomrule
\end{longtable}

\subsection{Supuestos y dependencias}

Se confirmo que Agricola Moreira cuenta con dispositivos moviles y computadores disponibles para el personal tecnico y administrativo. La colaboracion del personal de campo, inicialmente asumida como voluntaria, quedo confirmada con la realizacion de las 17 entrevistas y las 7 sesiones de validacion descritas en la Seccion 3.

\section{Requisitos especificos}

Esta seccion reune los 39 requisitos funcionales y los 21 requisitos no funcionales del sistema, obtenidos de las 17 entrevistas semiestructuradas y de las 7 sesiones de validacion con walkthrough, 3 con perfiles tecnicos y 4 con perfiles no tecnicos. De los 39 requisitos funcionales, 31 cuentan con evidencia directa de entrevista y 8 se derivaron directamente de los 26 criterios de cumplimiento legal descritos en \ruta{Modelo\_Legal\_LOPDP.md}, en sus tres bloques normativos.

Cada entrevista tiene asignado un codigo de evidencia. EV-01 corresponde a ENTR-01, administrador. EV-02 a ENTR-02, jornalero. EV-03 a ENTR-03, jornalera. EV-04 a ENTR-04, trabajador. EV-05 y EV-06 a ENTR-05 y ENTR-06, tecnicos. EV-07 a ENTR-07, tecnico, en sesion de walkthrough. EV-08 a ENTR-08, jornalera, en sesion de walkthrough. EV-09 a ENTR-09, jornalero. EV-10 a ENTR-10, tecnico, en sesion de walkthrough. EV-11 y EV-12 a ENTR-11 y ENTR-12, jornaleros, en sesion de walkthrough. EV-13 y EV-14 a ENTR-13 y ENTR-14, jornaleros. EV-15 a ENTR-15, jornalero, en sesion de walkthrough. EV-16 a ENTR-16, tecnica, en sesion de walkthrough. EV-17 a ENTR-17, tecnico. Las 17 personas entrevistadas son distintas entre si, aunque algunas comparten nombre de pila.

\subsection{Requisitos funcionales con evidencia de entrevista}

\rf{RF-01}{Registro de parcelas y lotes}
\begin{itemize}
\item \textbf{Descripcion:} el sistema debe permitir registrar y editar una parcela o lote con nombre o numero, sector, ubicacion, area, tipo de cultivo, variedad, cantidad de plantas y estado, activo o inactivo.
\item \textbf{Actor y origen:} Administrador, Tecnico, Jornalero. EV-01, EV-07, EV-09, EV-10, EV-11, EV-14, EV-15, EV-16.
\item \textbf{Entradas:} nombre, sector, ubicacion, area, cultivo, variedad, numero de plantas.
\item \textbf{Salidas:} ficha de parcela registrada.
\item \textbf{Precondicion:} usuario autenticado con permiso de edicion.
\item \textbf{Postcondicion:} parcela visible en el listado general.
\item \textbf{Flujo principal:}
\begin{enumerate}[label=\arabic*.,leftmargin=1.6em]
\item El usuario (Administrador, Tecnico o Jornalero) autenticado con permiso de edicion selecciona ``Nueva parcela''.
\item Ingresa nombre/numero, sector, ubicacion, area, cultivo, variedad y cantidad de plantas.
\item El sistema valida que los 7 campos obligatorios esten completos.
\item El sistema guarda la parcela y la muestra en el listado general en menos de 2 segundos.
\end{enumerate}
\item \textbf{Flujos alternativos y excepciones:}
\begin{itemize}
\item Si falta alguno de los 7 campos obligatorios, el sistema rechaza el guardado y marca los campos vacios (criterio de verificacion de RF-01).
\item Si la parcela se marca como ``para exportacion'', el sistema habilita los campos de trazabilidad de exportacion (ver RF-14, CU-03).
\end{itemize}
\item \textbf{Prioridad MoSCoW:} Must have.
\item \textbf{Criterio de verificacion:} al crear una parcela con los 7 campos obligatorios, el sistema la muestra en el listado en menos de 2 segundos, sin campos vacios.
\item \textbf{Evidencia textual:} ``Saber exactamente el lote por su nombre y numero, que cultivo tiene y cuantas plantas han sido o estan en tratamiento'' (EV-15). ``Los datos deberian ser el codigo, el tamano, el cultivo y la ubicacion'' (EV-16).
\end{itemize}

\rf{RF-02}{Catalogo cerrado de cultivos y variedades}
\begin{itemize}
\item \textbf{Descripcion:} el sistema debe ofrecer una lista predefinida y editable por el administrador de cultivos, cacao, platano, palma y ampliable, y sus variedades, sin permitir texto libre del tipo ``otros''.
\item \textbf{Actor y origen:} Jornalero, Tecnico. EV-10, EV-11, EV-12.
\item \textbf{Entradas:} seleccion de cultivo o variedad desde el catalogo.
\item \textbf{Salidas:} cultivo o variedad asociado a la parcela.
\item \textbf{Precondicion:} catalogo previamente configurado.
\item \textbf{Postcondicion:} ningun registro de cultivo queda con valor ``otros'' sin especificar.
\item \textbf{Flujo principal:}
\begin{enumerate}[label=\arabic*.,leftmargin=1.6em]
\item El usuario, al registrar o editar una parcela (CU-01), abre el selector de cultivo/variedad.
\item El sistema muestra unicamente las opciones del catalogo previamente configurado por el Administrador.
\item El usuario selecciona el cultivo/variedad y el sistema lo asocia a la parcela.
\end{enumerate}
\item \textbf{Flujos alternativos y excepciones:}
\begin{itemize}
\item Si el usuario intenta ingresar un cultivo fuera del catalogo (texto libre tipo ``otros''), el sistema rechaza el valor (criterio de verificacion de RF-02).
\item El Administrador puede ampliar el catalogo dando de alta un nuevo cultivo/variedad, lo que lo deja disponible de inmediato para el resto de usuarios.
\end{itemize}
\item \textbf{Prioridad MoSCoW:} Should have.
\item \textbf{Criterio de verificacion:} intentar guardar un cultivo fuera del catalogo es rechazado por el sistema.
\item \textbf{Evidencia textual:} ``no que este platano, cacao ahi... Simplemente que permite escribir... No deberia estar asi'' (EV-11). ``agrupadas por secciones, cacao, platano, etcetera'' (EV-12).
\item \textbf{Nota de relacion:} este RF se amplia en RF-34, que exige que el catalogo soporte ademas un conjunto de atributos distinto por variedad, no solo por cultivo.
\end{itemize}

\rf{RF-03}{Registro de actividades agricolas}
\begin{itemize}
\item \textbf{Descripcion:} el sistema debe permitir registrar una actividad, fumigacion, fertilizacion, poda, limpieza o riego, indicando fecha, lote, producto o insumo usado y trabajador responsable.
\item \textbf{Criterio legal cubierto:} C14 (Res. AGROCALIDAD 183, Art. 3).
\item \textbf{Actor y origen:} Jornalero, Tecnico. EV-03, EV-07, EV-08, EV-10, EV-14, EV-15, EV-16.
\item \textbf{Entradas:} fecha, lote, tipo de actividad, producto, trabajador.
\item \textbf{Salidas:} registro de actividad.
\item \textbf{Precondicion:} lote existente.
\item \textbf{Postcondicion:} actividad queda asociada al historial del lote.
\item \textbf{Flujo principal:}
\begin{enumerate}[label=\arabic*.,leftmargin=1.6em]
\item El Jornalero o Tecnico selecciona el lote existente sobre el que va a registrar la actividad.
\item Ingresa fecha, tipo de actividad (fumigacion, fertilizacion, poda, riego, limpieza), producto/insumo usado y trabajador responsable.
\item El sistema guarda el registro y lo asocia al historial del lote (CU-02).
\end{enumerate}
\item \textbf{Flujos alternativos y excepciones:}
\begin{itemize}
\item Si falta la fecha, el lote o el trabajador responsable, el sistema rechaza el guardado (criterio de verificacion de RF-03).
\item Si la actividad usa un insumo del catalogo, el sistema descuenta automaticamente el inventario correspondiente (ver RF-07, CU-05).
\end{itemize}
\item \textbf{Prioridad MoSCoW:} Must have.
\item \textbf{Criterio de verificacion:} el sistema no permite guardar una actividad sin fecha, lote y trabajador.
\item \textbf{Evidencia textual:} ``que se registre lo que es la fecha, la parcela, la actividad, producto y el trabajador'' (EV-10). ``La informacion dispensable seria la actividad, la parcela, la fecha y tambien el responsable'' (EV-16).
\end{itemize}

\rf{RF-04}{Registro de cosecha por unidad especifica del cultivo}
\begin{itemize}
\item \textbf{Descripcion:} el sistema debe registrar la cosecha usando la unidad propia de cada cultivo, racimo o caja para platano, tacho o quintal para cacao, con conversion configurable a kilogramos o libras.
\item \textbf{Actor y origen:} Jornalero. EV-02, EV-08, EV-11, EV-12, EV-13, EV-14, EV-15.
\item \textbf{Entradas:} lote, fecha, cultivo, cantidad, unidad.
\item \textbf{Salidas:} registro de cosecha.
\item \textbf{Precondicion:} lote y cultivo existentes.
\item \textbf{Postcondicion:} cosecha sumada al acumulado del lote o periodo.
\item \textbf{Flujo principal:}
\begin{enumerate}[label=\arabic*.,leftmargin=1.6em]
\item El Jornalero selecciona el lote y el cultivo sobre el que va a registrar la cosecha.
\item El sistema determina la unidad esperada segun el cultivo (tacho/quintal para cacao, racimo/caja para platano).
\item El usuario ingresa la cantidad cosechada en esa unidad.
\item El sistema valida la unidad, guarda la cosecha y actualiza el acumulado del lote/periodo.
\end{enumerate}
\item \textbf{Flujos alternativos y excepciones:}
\begin{itemize}
\item Si la unidad ingresada no corresponde al cultivo del lote, el sistema rechaza el registro (criterio de verificacion de RF-04).
\item Si corresponde, el usuario puede consultar de inmediato el precio de mercado vigente y ver el ingreso proyectado (RF-15).
\end{itemize}
\item \textbf{Prioridad MoSCoW:} Must have.
\item \textbf{Criterio de verificacion:} al registrar cosecha de cacao el sistema exige unidad tacho o quintal, y de platano racimo o caja.
\item \textbf{Evidencia textual:} ``cuando es la cosecha de platano, es por racimo. Y cuando es el cacao, por tacho'' (EV-02). ``Me gustaria el total de la produccion y tambien lo que es del peso'' (EV-08).
\end{itemize}

\rf{RF-05}{Calculo automatico de rendimiento por lote y periodo}
\begin{itemize}
\item \textbf{Descripcion:} el sistema debe calcular automaticamente el promedio de produccion por lote y periodo, semanal o mensual, y permitir comparar entre periodos.
\item \textbf{Actor y origen:} Administrador, Jornalero, Tecnica. EV-01, EV-07, EV-11, EV-15, EV-16.
\item \textbf{Entradas:} historico de cosechas del lote.
\item \textbf{Salidas:} promedio y comparacion entre periodos.
\item \textbf{Precondicion:} existen 2 o mas registros de cosecha del lote.
\item \textbf{Postcondicion:} valor mostrado en el reporte del lote.
\item \textbf{Flujo principal:}
\begin{enumerate}[label=\arabic*.,leftmargin=1.6em]
\item El sistema consulta el historico de cosechas del lote seleccionado.
\item Calcula el promedio de produccion del periodo (suma dividida entre N registros).
\item Muestra el promedio y permite comparar contra otros periodos en el reporte del lote.
\end{enumerate}
\item \textbf{Flujos alternativos y excepciones:}
\begin{itemize}
\item Si el lote tiene menos de 2 registros de cosecha, el sistema muestra ``datos insuficientes'' en vez de un promedio (precondicion de RF-05).
\end{itemize}
\item \textbf{Prioridad MoSCoW:} Must have.
\item \textbf{Criterio de verificacion:} el promedio calculado coincide con el calculo manual, suma dividida entre N, sobre un conjunto de prueba.
\item \textbf{Evidencia textual:} ``sacar un promedio de cuanto rinde cada lote'' (EV-15). ``Me gustaria que generara automaticamente lo que es el rendimiento y tambien los totales'' (EV-16).
\end{itemize}

\rf{RF-06}{Calculo automatico de ganancia neta}
\begin{itemize}
\item \textbf{Descripcion:} el sistema debe calcular ingresos menos egresos, insumos y mano de obra, y mostrar la ganancia neta por periodo.
\item \textbf{Actor y origen:} Tecnico, Jornalero. EV-10, EV-11.
\item \textbf{Entradas:} registros de ingresos y egresos.
\item \textbf{Salidas:} ganancia neta del periodo.
\item \textbf{Precondicion:} existen registros de ingreso y egreso en el periodo.
\item \textbf{Postcondicion:} valor visible en el modulo financiero.
\item \textbf{Flujo principal:}
\begin{enumerate}[label=\arabic*.,leftmargin=1.6em]
\item El sistema consulta los registros de ingreso (automaticos por cosecha/precio y manuales por RF-26) y egreso (insumos, mano de obra) del periodo.
\item Calcula la ganancia neta como la suma de ingresos menos la suma de egresos.
\item Muestra el valor en el modulo financiero.
\end{enumerate}
\item \textbf{Flujos alternativos y excepciones:}
\begin{itemize}
\item Si no existen registros de ingreso o egreso en el periodo consultado, el sistema muestra el modulo financiero en cero, sin error.
\item El usuario puede registrar manualmente un ingreso o egreso adicional (RF-26), que se refleja de inmediato en el siguiente calculo.
\end{itemize}
\item \textbf{Prioridad MoSCoW:} Must have.
\item \textbf{Criterio de verificacion:} la ganancia es igual a la suma de ingresos menos la suma de egresos, verificado contra un caso de prueba manual.
\item \textbf{Evidencia textual:} ``de todo lo que se gasto se ingresa tambien todo lo que se vendio y menorando todo eso, lo que quedo de ganancia'' (EV-10).
\item \textbf{Nota de relacion:} ver tambien RF-26, que pide un modulo de registro manual de ingresos y egresos independiente de este calculo automatico.
\end{itemize}

\rf{RF-07}{Gestion de inventario de insumos}
\begin{itemize}
\item \textbf{Descripcion:} el sistema debe permitir dar de alta insumos, registrar cantidad disponible o usada y visualizar el stock restante.
\item \textbf{Actor y origen:} Administrador, Tecnico, Jornalero. EV-01, EV-03, EV-05, EV-06, EV-09, EV-10, EV-11, EV-12, EV-16.
\item \textbf{Entradas:} nombre del insumo, cantidad, unidad.
\item \textbf{Salidas:} stock disponible actualizado.
\item \textbf{Precondicion:} ninguna.
\item \textbf{Postcondicion:} stock reducido tras cada uso registrado.
\item \textbf{Flujo principal:}
\begin{enumerate}[label=\arabic*.,leftmargin=1.6em]
\item El usuario da de alta un insumo con nombre, cantidad disponible y unidad, o registra el uso de una cantidad de un insumo existente.
\item El sistema actualiza el stock disponible restando la cantidad usada.
\item El sistema muestra el stock actualizado.
\end{enumerate}
\item \textbf{Flujos alternativos y excepciones:}
\begin{itemize}
\item Si el uso registrado dejaria el stock en un valor negativo, el sistema advierte antes de permitir guardar (regla de negocio de CU-05).
\item Tras actualizar el stock, el sistema verifica si quedo por debajo del umbral configurado y dispara la alerta de bajo stock (RF-08, RF-30).
\end{itemize}
\item \textbf{Prioridad MoSCoW:} Must have.
\item \textbf{Criterio de verificacion:} tras registrar el uso de 26 unidades de un stock de 30, el sistema muestra 4 disponibles.
\item \textbf{Evidencia textual:} ``dice 4 de 30, es decir, lo que tenia era 30 y ahora le queda 4'' (EV-10). ``Sobre cada insumo, lo que es la cantidad, el nombre y ademas del stock limit'' (EV-16).
\end{itemize}

\rf{RF-08}{Alertas automaticas de bajo stock}
\begin{itemize}
\item \textbf{Descripcion:} el sistema debe enviar una alerta, notificacion, SMS o correo, cuando el stock de un insumo baje de un umbral configurable, anticipandose al calendario de fertilizacion o fumigacion.
\item \textbf{Actor y origen:} Administrador, Jornalero, Tecnica. EV-01, EV-03, EV-08, EV-09, EV-10, EV-11, EV-16.
\item \textbf{Entradas:} umbral configurado, stock actual.
\item \textbf{Salidas:} notificacion al responsable.
\item \textbf{Precondicion:} umbral definido para el insumo.
\item \textbf{Postcondicion:} notificacion enviada y registrada.
\item \textbf{Flujo principal:}
\begin{enumerate}[label=\arabic*.,leftmargin=1.6em]
\item Tras cada actualizacion de stock de un insumo (RF-07), el sistema compara el nuevo stock contra el umbral minimo configurado (RF-30).
\item Si el stock quedo por debajo del umbral, el sistema genera una notificacion al responsable del insumo.
\item El sistema registra el envio de la notificacion.
\end{enumerate}
\item \textbf{Flujos alternativos y excepciones:}
\begin{itemize}
\item Si el insumo no tiene un umbral configurado todavia, el sistema no dispara la alerta hasta que el Tecnico lo defina (RF-30).
\item Si la notificacion no puede entregarse por falta de conectividad, el sistema reintenta hasta la reconexion del dispositivo (ver RNF-15).
\end{itemize}
\item \textbf{Prioridad MoSCoW:} Must have.
\item \textbf{Criterio de verificacion:} al bajar del umbral, la notificacion se envia en menos de 5 minutos.
\item \textbf{Evidencia textual:} ``en siete dias esta estimado que debes abonar el cacao... Y se necesita comprar'' (EV-11). ``Me gustaria que me avisara por medio de una notificacion en el celular'' (EV-08).
\item \textbf{Nota de relacion:} ver RF-30 para el umbral configurable como campo propio del insumo.
\end{itemize}

\rf{RF-09}{Alerta de plagas o enfermedades asistida por IA con verificacion humana obligatoria}
\begin{itemize}
\item \textbf{Descripcion:} el sistema debe generar una alerta de posible plaga o enfermedad por lote, basada en los datos ingresados, marcada explicitamente como sugerencia a confirmar en campo, nunca como diagnostico definitivo automatico.
\item \textbf{Actor y origen:} Administrador, Tecnico, Jornalero. EV-01, EV-06, EV-07, EV-08, EV-09, EV-10, EV-15, EV-16, EV-17.
\item \textbf{Entradas:} datos del lote, historico de tratamientos.
\item \textbf{Salidas:} alerta con opcion de confirmar o descartar.
\item \textbf{Precondicion:} modelo de IA entrenado disponible.
\item \textbf{Postcondicion:} alerta queda registrada con el resultado de la verificacion humana.
\item \textbf{Flujo principal:}
\begin{enumerate}[label=\arabic*.,leftmargin=1.6em]
\item El sistema de IA analiza los datos del lote (historico de tratamientos y, opcionalmente, una foto) y genera una sugerencia de posible plaga/enfermedad con su justificacion.
\item El sistema notifica al usuario responsable, marcando explicitamente la sugerencia como ``a confirmar en campo'', nunca como diagnostico automatico.
\item El usuario revisa la justificacion y elige ``confirmar'' o ``descartar''.
\item El sistema registra la decision del usuario junto con la alerta.
\end{enumerate}
\item \textbf{Flujos alternativos y excepciones:}
\begin{itemize}
\item El usuario descarta la sugerencia; el sistema la deja registrada como descartada, sin aplicar ninguna accion sobre el lote.
\item El usuario, en desacuerdo con la recomendacion, la registra formalmente en la bandeja de sugerencias en vez de solo confirmar/descartar (RF-33).
\item En ningun caso la alerta se aplica automaticamente al lote sin la confirmacion explicita del usuario responsable (regla de negocio central de RF-09).
\end{itemize}
\item \textbf{Prioridad MoSCoW:} Must have.
\item \textbf{Criterio de verificacion:} ninguna alerta se aplica automaticamente a un lote sin que el usuario responsable la confirme.
\item \textbf{Evidencia textual:} ``yo confirmaria yendo a visualizar el terreno'' (EV-10). ``Se deberia poder revisar antes de tomar una decision o aplicacion'' (EV-07, EV-08). ``Primeramente que eso venga de una fuente confiable o de ya experiencias realizadas'' (EV-17).
\end{itemize}

\rf{RF-10}{Diagnostico de plagas por imagen}
\begin{itemize}
\item \textbf{Descripcion:} el sistema debe permitir subir una foto de una hoja o fruto y recibir una sugerencia de plaga o enfermedad probable y tratamiento recomendado.
\item \textbf{Actor y origen:} Jornalero. EV-15.
\item \textbf{Entradas:} foto en JPG o PNG.
\item \textbf{Salidas:} plaga probable y recomendacion.
\item \textbf{Precondicion:} conexion a internet.
\item \textbf{Postcondicion:} resultado registrado en el historial del lote.
\item \textbf{Flujo principal:}
\begin{enumerate}[label=\arabic*.,leftmargin=1.6em]
\item El Jornalero sube una foto (JPG/PNG) de una hoja o fruto desde el modulo de alerta de plagas (CU-06).
\item El sistema procesa la imagen y genera una sugerencia de plaga/enfermedad probable junto con el tratamiento recomendado.
\item El sistema responde en menos de 10 segundos y registra el resultado en el historial del lote.
\end{enumerate}
\item \textbf{Flujos alternativos y excepciones:}
\begin{itemize}
\item Si no hay conexion a internet, el sistema informa que el diagnostico por imagen no esta disponible en ese momento (precondicion de RF-10).
\item El resultado del diagnostico por imagen sigue el mismo flujo de confirmacion humana obligatoria que RF-09: se muestra como sugerencia, nunca como diagnostico definitivo.
\end{itemize}
\item \textbf{Prioridad MoSCoW:} Could have.
\item \textbf{Criterio de verificacion:} el sistema responde con una sugerencia en menos de 10 segundos tras subir la foto.
\item \textbf{Evidencia textual:} ``una foto de una hoja o de un fruto enfermo que el sistema le atiene rapido y que plaga es, que remedio exacto'' (EV-15).
\end{itemize}

\rf{RF-11}{Asignacion y seguimiento de tareas}
\begin{itemize}
\item \textbf{Descripcion:} el sistema debe permitir asignar una tarea a un trabajador por lote y fecha, con estados pendiente, en proceso, bloqueado y completado, y filtro por estado.
\item \textbf{Actor y origen:} Administrador, Tecnico, Jornalero. EV-01, EV-09, EV-10, EV-12, EV-16.
\item \textbf{Entradas:} lote, trabajador, fecha, tipo de tarea.
\item \textbf{Salidas:} tarea visible con su estado.
\item \textbf{Precondicion:} trabajador registrado.
\item \textbf{Postcondicion:} tarea filtrable por los 4 estados.
\item \textbf{Flujo principal:}
\begin{enumerate}[label=\arabic*.,leftmargin=1.6em]
\item El usuario asigna una tarea a un trabajador, indicando lote, fecha y tipo de tarea.
\item El sistema crea la tarea en estado ``pendiente'' y la muestra con su estado.
\item El trabajador actualiza el estado de la tarea (pendiente → en proceso → bloqueado → completado) conforme avanza.
\item El sistema permite filtrar la lista de tareas por cualquiera de los 4 estados.
\end{enumerate}
\item \textbf{Flujos alternativos y excepciones:}
\begin{itemize}
\item Si el lote tiene un riesgo laboral registrado (CU-11), el sistema muestra la advertencia de equipo de proteccion antes de confirmar la asignacion.
\item Al filtrar por un estado, el sistema muestra exclusivamente las tareas en ese estado, sin mezclar lotes ni trabajadores (criterio de verificacion de RF-11).
\end{itemize}
\item \textbf{Prioridad MoSCoW:} Must have.
\item \textbf{Criterio de verificacion:} al filtrar por el estado pendiente, el sistema solo muestra tareas en ese estado.
\item \textbf{Evidencia textual:} ``que tengas esa opcion de buscar en esas cuatro secciones'' (EV-12).
\end{itemize}

\rf{RF-12}{Notificacion movil de tareas asignadas}
\begin{itemize}
\item \textbf{Descripcion:} el sistema debe notificar al trabajador, por push o SMS, el lote, la fecha y la tarea asignada.
\item \textbf{Actor y origen:} Tecnico, Jornalera. EV-07, EV-08, EV-10.
\item \textbf{Entradas:} asignacion de tarea.
\item \textbf{Salidas:} notificacion push o SMS.
\item \textbf{Precondicion:} trabajador con dispositivo registrado.
\item \textbf{Postcondicion:} notificacion entregada.
\item \textbf{Flujo principal:}
\begin{enumerate}[label=\arabic*.,leftmargin=1.6em]
\item Al asignarse una tarea (RF-11), el sistema identifica el dispositivo registrado del trabajador.
\item El sistema envia una notificacion push/SMS con el lote, la fecha y el tipo de tarea.
\item El trabajador recibe la notificacion en su dispositivo.
\end{enumerate}
\item \textbf{Flujos alternativos y excepciones:}
\begin{itemize}
\item Si el trabajador no tiene un dispositivo registrado, el sistema deja la tarea visible en su vista de pendientes (RF-27) aunque no pueda enviar la notificacion push.
\end{itemize}
\item \textbf{Prioridad MoSCoW:} Should have.
\item \textbf{Criterio de verificacion:} la notificacion incluye lote, fecha y tipo de tarea en el mensaje.
\item \textbf{Evidencia textual:} ``que la notificacion... Le avise por lote, en tal lote, fecha tal, hay un control'' (EV-10).
\end{itemize}

\rf{RF-13}{Registro de motivo de rechazo o perdida de fruta}
\begin{itemize}
\item \textbf{Descripcion:} el sistema debe permitir registrar, por lote y fecha, la cantidad de fruta rechazada y su motivo, enfermedad, plaga o dano mecanico.
\item \textbf{Actor y origen:} Jornalero. EV-14.
\item \textbf{Entradas:} lote, fecha, cantidad rechazada, motivo.
\item \textbf{Salidas:} registro de rechazo.
\item \textbf{Precondicion:} cosecha registrada ese dia.
\item \textbf{Postcondicion:} rechazo restado del total aprovechable.
\item \textbf{Flujo principal:}
\begin{enumerate}[label=\arabic*.,leftmargin=1.6em]
\item El Jornalero, tras registrar la cosecha del dia (RF-04), registra la cantidad de fruta rechazada y su motivo (enfermedad, plaga, dano mecanico).
\item El sistema resta automaticamente la cantidad rechazada del total aprovechable del lote.
\end{enumerate}
\item \textbf{Flujos alternativos y excepciones:}
\begin{itemize}
\item Si no hay una cosecha registrada ese dia, el sistema no permite registrar un rechazo asociado (precondicion de RF-13).
\end{itemize}
\item \textbf{Prioridad MoSCoW:} Should have.
\item \textbf{Criterio de verificacion:} el total de fruta aprovechable excluye automaticamente lo marcado como rechazado.
\item \textbf{Evidencia textual:} ``cuantos salieron rechazados, ya sea por una enfermedad, una plaga o los danos'' (EV-14).
\end{itemize}

\rf{RF-14}{Trazabilidad de lote de exportacion}
\begin{itemize}
\item \textbf{Descripcion:} el sistema debe permitir registrar, por lote de exportacion, el estado de calidad, primera o segunda, la variedad sembrada y el ciclo de enfunde, fecha y color de cinta, para trazar el producto desde el campo hasta el empaque.
\item \textbf{Criterio legal cubierto:} C15 (Res. AGROCALIDAD 183, Art. 36).
\item \textbf{Actor y origen:} Jornalero. EV-13, EV-14.
\item \textbf{Entradas:} color de cinta, semana de enfunde, calidad de caja.
\item \textbf{Salidas:} ficha de trazabilidad del lote de exportacion.
\item \textbf{Precondicion:} lote marcado como exportacion.
\item \textbf{Postcondicion:} historial completo desde el enfunde hasta el empaque consultable.
\item \textbf{Flujo principal:}
\begin{enumerate}[label=\arabic*.,leftmargin=1.6em]
\item Para un lote marcado como ``exportacion'' (RF-01), el usuario registra color de cinta, semana de enfunde y calidad de caja al momento de la cosecha (RF-04).
\item El sistema guarda la ficha de trazabilidad asociada al lote de exportacion.
\item Dado un numero de caja, el sistema recupera el lote, la fecha de enfunde y el color de cinta correspondiente (criterio de verificacion de RF-14).
\end{enumerate}
\item \textbf{Flujos alternativos y excepciones:}
\begin{itemize}
\item Si el lote no esta marcado como ``exportacion'', el sistema no solicita estos campos adicionales.
\end{itemize}
\item \textbf{Prioridad MoSCoW:} Should have.
\item \textbf{Criterio de verificacion:} dado un numero de caja, el sistema recupera el lote, la fecha de enfunde y el color de cinta correspondiente.
\item \textbf{Evidencia textual:} ``con el dato que el enfundador da, con eso se lo maneja para el dia de cosecha'' (EV-13). ``que en las pantallas de la cosecha se debe ver clarito el color de la cinta'' (EV-14).
\end{itemize}

\rf{RF-15}{Consulta de precio de mercado y estimacion de ingreso}
\begin{itemize}
\item \textbf{Descripcion:} el sistema debe permitir consultar el precio de mercado vigente del cacao o del platano y estimar el ingreso proyectado de una cosecha.
\item \textbf{Actor y origen:} Jornalero. EV-11.
\item \textbf{Entradas:} cantidad cosechada, precio de mercado, manual o de fuente externa.
\item \textbf{Salidas:} ingreso estimado.
\item \textbf{Precondicion:} cantidad de cosecha registrada.
\item \textbf{Postcondicion:} estimacion mostrada al usuario.
\item \textbf{Flujo principal:}
\begin{enumerate}[label=\arabic*.,leftmargin=1.6em]
\item Tras registrar una cosecha (RF-04), el usuario consulta el precio de mercado vigente del cultivo.
\item El sistema calcula el ingreso estimado como la cantidad cosechada multiplicada por el precio ingresado.
\item El sistema muestra la estimacion al usuario.
\end{enumerate}
\item \textbf{Flujos alternativos y excepciones:}
\begin{itemize}
\item Si no hay un precio de mercado disponible de fuente externa, el usuario puede ingresarlo manualmente para obtener la estimacion.
\end{itemize}
\item \textbf{Prioridad MoSCoW:} Could have.
\item \textbf{Criterio de verificacion:} el ingreso estimado es igual a la cantidad multiplicada por el precio ingresado, validado con un caso de prueba.
\item \textbf{Evidencia textual:} ``que consultara el precio de la materia prima... Un calculo de cuanto mas o menos uno va a ganar'' (EV-11).
\end{itemize}

\rf{RF-16}{Registro de condicion climatica y estado del lote}
\begin{itemize}
\item \textbf{Descripcion:} el sistema debe permitir anotar, junto a cada visita o actividad del lote, el estado climatico, lluvia o sol, y del terreno, accesibilidad.
\item \textbf{Actor y origen:} Jornalero. EV-12, EV-15.
\item \textbf{Entradas:} condicion climatica, observacion de acceso.
\item \textbf{Salidas:} campo adicional en el registro de actividad o parcela.
\item \textbf{Precondicion:} ninguna.
\item \textbf{Postcondicion:} dato consultable en el historial del lote.
\item \textbf{Flujo principal:}
\begin{enumerate}[label=\arabic*.,leftmargin=1.6em]
\item Al registrar una actividad o visitar un lote (CU-01), el usuario anota la condicion climatica (lluvia/sol) y una observacion de accesibilidad del terreno.
\item El sistema guarda el dato como campo adicional del registro de actividad/parcela.
\end{enumerate}
\item \textbf{Flujos alternativos y excepciones:}
\begin{itemize}
\item El usuario puede registrar la condicion climatica en cualquier momento posterior, no solo al crear el registro (flujo alternativo documentado en CU-01).
\end{itemize}
\item \textbf{Prioridad MoSCoW:} Could have.
\item \textbf{Criterio de verificacion:} el campo de clima aparece disponible en el formulario de registro de parcela o actividad.
\item \textbf{Evidencia textual:} ``Puede faltar un espacio para anotar lo que es el estado del clima'' (EV-15).
\end{itemize}

\rf{RF-17}{Canal de reporte rapido por chat o voz}
\begin{itemize}
\item \textbf{Descripcion:} el sistema debe ofrecer un boton de reporte rapido, texto corto o nota de voz, para que un jornalero reporte una novedad sin llenar formularios extensos.
\item \textbf{Actor y origen:} Jornalero. EV-15.
\item \textbf{Entradas:} nota de voz o texto corto.
\item \textbf{Salidas:} reporte enviado al responsable del lote.
\item \textbf{Precondicion:} usuario autenticado.
\item \textbf{Postcondicion:} reporte visible para el administrador o tecnico.
\item \textbf{Flujo principal:}
\begin{enumerate}[label=\arabic*.,leftmargin=1.6em]
\item El Jornalero abre el canal de reporte rapido desde la pantalla principal.
\item Registra un mensaje corto (texto o nota de voz) describiendo una novedad de campo.
\item El sistema envia el reporte al responsable del lote en un maximo de 2 toques desde la pantalla principal (criterio de verificacion de RF-17).
\end{enumerate}
\item \textbf{Flujos alternativos y excepciones:}
\begin{itemize}
\item Si el reporte describe una posible plaga cuarentenaria, el sistema sugiere escalarlo como aviso fitosanitario formal (RF-37, CU-14).
\end{itemize}
\item \textbf{Prioridad MoSCoW:} Could have.
\item \textbf{Criterio de verificacion:} el reporte se envia en un maximo de 2 toques desde la pantalla principal.
\item \textbf{Evidencia textual:} ``un chat o un boton rapido de voz para reportar novedades'' (EV-15).
\end{itemize}

\rf{RF-18}{Registro de riesgo y seguridad laboral por lote}
\begin{itemize}
\item \textbf{Descripcion:} el sistema debe permitir marcar un lote con una alerta de riesgo, fauna peligrosa o terreno inestable, visible antes de asignar tareas ahi, y sugerir el equipo de proteccion personal requerido.
\item \textbf{Criterio legal cubierto:} C16 (Res. AGROCALIDAD 183, Art. 38); C22 (Res. AGROCALIDAD 0072, Art. 18(g) Res.183).
\item \textbf{Actor y origen:} Jornalero, Tecnico. EV-11, EV-17.
\item \textbf{Entradas:} tipo de riesgo, lote, equipo de proteccion recomendado.
\item \textbf{Salidas:} alerta visible al asignar tarea en ese lote.
\item \textbf{Precondicion:} ninguna.
\item \textbf{Postcondicion:} advertencia mostrada al asignar personal.
\item \textbf{Flujo principal:}
\begin{enumerate}[label=\arabic*.,leftmargin=1.6em]
\item El Tecnico marca un lote con un riesgo laboral (fauna peligrosa, terreno inestable), indicando el equipo de proteccion personal sugerido.
\item Al asignar una tarea en ese lote (RF-11), el sistema muestra la advertencia de riesgo antes de confirmar la asignacion.
\end{enumerate}
\item \textbf{Flujos alternativos y excepciones:}
\begin{itemize}
\item Ninguna tarea sobre un lote con riesgo registrado puede confirmarse sin que la advertencia de equipo de proteccion se haya mostrado primero (regla de negocio de CU-11).
\end{itemize}
\item \textbf{Prioridad MoSCoW:} Should have.
\item \textbf{Criterio de verificacion:} al asignar una tarea en un lote marcado con riesgo, el sistema muestra la advertencia antes de confirmar.
\item \textbf{Evidencia textual:} ``hay riesgo de que haya animales... Como serpientes'' (EV-11). ``se dice que usar, que trajes usar, incluso botas... Ya que al momento de presenciar estos peligros esten bien cuidados'' (EV-17).
\end{itemize}

\rf{RF-19}{Reportes de produccion, costos y perdidas con graficos}
\begin{itemize}
\item \textbf{Descripcion:} el sistema debe generar reportes de produccion, costos, perdidas por plaga y rendimiento por parcela o periodo, con graficos de barra y circulares codificados por color.
\item \textbf{Criterio legal cubierto:} C18 (Res. AGROCALIDAD 183, Art. 27); C19 (Res. AGROCALIDAD 183, Art. 29).
\item \textbf{Actor y origen:} Administrador, Tecnico. EV-01, EV-05, EV-06, EV-07, EV-10, EV-16.
\item \textbf{Entradas:} rango de fechas, lote o lotes.
\item \textbf{Salidas:} reporte visual exportable.
\item \textbf{Precondicion:} existen datos en el rango seleccionado.
\item \textbf{Postcondicion:} reporte generado y descargable.
\item \textbf{Flujo principal:}
\begin{enumerate}[label=\arabic*.,leftmargin=1.6em]
\item El Administrador o Tecnico selecciona un rango de fechas y uno o mas lotes.
\item El sistema consolida produccion, costos, perdidas por plaga y rendimiento del rango seleccionado.
\item El sistema genera graficos de barra y circulares codificados por color, reflejando exactamente la suma de los registros del rango.
\end{enumerate}
\item \textbf{Flujos alternativos y excepciones:}
\begin{itemize}
\item Si no existen datos en el rango seleccionado, el sistema muestra ``no hay datos disponibles'' en vez de un reporte vacio (criterio de verificacion de RF-19).
\end{itemize}
\item \textbf{Prioridad MoSCoW:} Must have.
\item \textbf{Criterio de verificacion:} el reporte generado refleja exactamente la suma de los registros del rango de fechas seleccionado.
\item \textbf{Evidencia textual:} ``reporte por parcela, costos, perdidas por plagas o rendimiento por periodo'' (EV-01).
\end{itemize}

\rf{RF-20}{Gestion de usuarios y control de acceso}
\begin{itemize}
\item \textbf{Descripcion:} el sistema debe requerir usuario y contrasena para ingresar, con roles diferenciados, administrador, tecnico y jornalero.
\item \textbf{Actor y origen:} Tecnico. EV-07, EV-08, EV-10.
\item \textbf{Entradas:} usuario, contrasena.
\item \textbf{Salidas:} sesion autenticada con permisos segun rol.
\item \textbf{Precondicion:} usuario previamente registrado.
\item \textbf{Postcondicion:} acceso concedido solo a las funciones de su rol.
\item \textbf{Flujo principal:}
\begin{enumerate}[label=\arabic*.,leftmargin=1.6em]
\item El usuario ingresa su usuario y contrasena.
\item El sistema valida las credenciales y determina el rol (Administrador, Tecnico, Jornalero).
\item El sistema concede acceso unicamente a las funciones correspondientes a ese rol.
\end{enumerate}
\item \textbf{Flujos alternativos y excepciones:}
\begin{itemize}
\item Si las credenciales son invalidas, el sistema niega el acceso.
\item Si un usuario con rol ``jornalero'' intenta acceder al modulo financiero, el sistema le niega el acceso (criterio de verificacion de RF-20).
\end{itemize}
\item \textbf{Prioridad MoSCoW:} Must have.
\item \textbf{Criterio de verificacion:} un usuario con rol jornalero no puede acceder al modulo financiero.
\item \textbf{Evidencia textual:} ``han creado como un tipo formulario, donde le pide lo que es un usuario y una contrasena'' (EV-10). ``El administrador deberia tener el control total y los trabajadores solo lo necesario'' (EV-07, EV-08).
\end{itemize}

\rf{RF-21}{Registro de visitas tecnicas periodicas}
\begin{itemize}
\item \textbf{Descripcion:} el sistema debe permitir registrar cada visita tecnica, fecha, hallazgos e informe adjunto, con recordatorio de la siguiente visita en un ciclo de 15 dias.
\item \textbf{Actor y origen:} Tecnico. EV-10.
\item \textbf{Entradas:} fecha de visita, informe, proxima fecha estimada.
\item \textbf{Salidas:} historial de visitas tecnicas por lote o finca.
\item \textbf{Precondicion:} ninguna.
\item \textbf{Postcondicion:} recordatorio generado para la siguiente visita.
\item \textbf{Flujo principal:}
\begin{enumerate}[label=\arabic*.,leftmargin=1.6em]
\item El Tecnico registra una visita tecnica: fecha, hallazgos e informe adjunto.
\item El sistema guarda la visita y la asocia al historial de cumplimiento del lote.
\item 15 dias despues, el sistema genera automaticamente un recordatorio de la siguiente visita (criterio de verificacion de RF-21).
\end{enumerate}
\item \textbf{Flujos alternativos y excepciones:}
\begin{itemize}
\item Si la visita encuentra un incumplimiento, el sistema sugiere registrar la capacitacion o el equipo de proteccion relacionado (CU-11, CU-13).
\end{itemize}
\item \textbf{Prioridad MoSCoW:} Should have.
\item \textbf{Criterio de verificacion:} 15 dias despues de una visita registrada, el sistema genera un recordatorio automatico.
\item \textbf{Evidencia textual:} ``un control de la visita tecnica de 15 dias y pone el otro control que fue a realizarla despues de 15 dias'' (EV-10).
\end{itemize}

\rf{RF-26}{Registro manual de ingresos y egresos}
\begin{itemize}
\item \textbf{Descripcion:} el sistema debe ofrecer un modulo dedicado para registrar manualmente cada ingreso y egreso, concepto, monto y fecha, independiente del calculo automatico de ganancia neta del RF-06.
\item \textbf{Actor y origen:} Tecnico, Jornalera. EV-07, EV-08.
\item \textbf{Entradas:} concepto, monto, fecha, tipo, ingreso o egreso.
\item \textbf{Salidas:} listado de movimientos.
\item \textbf{Precondicion:} usuario autenticado.
\item \textbf{Postcondicion:} movimiento reflejado en el calculo de ganancia neta.
\item \textbf{Flujo principal:}
\begin{enumerate}[label=\arabic*.,leftmargin=1.6em]
\item El usuario autenticado abre el modulo de registro manual de ingresos/egresos.
\item Registra concepto, monto, fecha y tipo (ingreso/egreso), de forma independiente al calculo automatico de RF-06.
\item El sistema guarda el movimiento en el listado.
\end{enumerate}
\item \textbf{Flujos alternativos y excepciones:}
\begin{itemize}
\item El movimiento registrado manualmente se refleja de inmediato en el reporte financiero del periodo correspondiente (criterio de verificacion de RF-26).
\end{itemize}
\item \textbf{Prioridad MoSCoW:} Must have.
\item \textbf{Criterio de verificacion:} un ingreso o egreso registrado manualmente aparece reflejado en el reporte financiero del periodo correspondiente.
\item \textbf{Evidencia textual:} ``Que permitiera controlar los ingresos y egresos'' (EV-07). ``Me gustaria que tuviera un registro de gastos y de ganancias tambien'' (EV-08).
\end{itemize}

\rf{RF-27}{Vista de tareas pendientes por trabajador}
\begin{itemize}
\item \textbf{Descripcion:} el sistema debe mostrar a cada trabajador una vista dedicada de sus propias tareas pendientes.
\item \textbf{Actor y origen:} Jornalero. EV-09.
\item \textbf{Entradas:} usuario autenticado.
\item \textbf{Salidas:} listado de tareas pendientes propias.
\item \textbf{Precondicion:} existen tareas asignadas al usuario.
\item \textbf{Postcondicion:} ninguna.
\item \textbf{Flujo principal:}
\begin{enumerate}[label=\arabic*.,leftmargin=1.6em]
\item El trabajador autenticado abre su vista de tareas pendientes.
\item El sistema muestra unicamente las tareas asignadas a ese trabajador.
\end{enumerate}
\item \textbf{Flujos alternativos y excepciones:}
\begin{itemize}
\item Un trabajador con 3 tareas pendientes ve exactamente esas 3 al abrir la vista, sin tareas de otros lotes o trabajadores (criterio de verificacion de RF-27).
\end{itemize}
\item \textbf{Prioridad MoSCoW:} Should have.
\item \textbf{Criterio de verificacion:} un trabajador con 3 tareas pendientes ve exactamente esas 3 al abrir la vista, sin tareas de otros lotes o trabajadores.
\item \textbf{Evidencia textual:} ``Existe alguna funcion que considere indispensable y que todavia no hayamos completado. Ver las tareas pendientes'' (EV-09).
\end{itemize}

\rf{RF-28}{Campo de observaciones en el registro de tareas}
\begin{itemize}
\item \textbf{Descripcion:} el sistema debe permitir adjuntar una nota u observacion libre a cada tarea registrada como completada.
\item \textbf{Actor y origen:} Jornalero. EV-09.
\item \textbf{Entradas:} texto libre.
\item \textbf{Salidas:} observacion asociada a la tarea.
\item \textbf{Precondicion:} tarea existente.
\item \textbf{Postcondicion:} observacion visible al consultar el historial de la tarea.
\item \textbf{Flujo principal:}
\begin{enumerate}[label=\arabic*.,leftmargin=1.6em]
\item El trabajador, al marcar una tarea como completada (RF-11), adjunta una nota u observacion de texto libre.
\item El sistema asocia la observacion a la tarea.
\end{enumerate}
\item \textbf{Flujos alternativos y excepciones:}
\begin{itemize}
\item La observacion ingresada permanece visible al consultar el detalle de la tarea en cualquier momento posterior (criterio de verificacion de RF-28).
\end{itemize}
\item \textbf{Prioridad MoSCoW:} Could have.
\item \textbf{Criterio de verificacion:} la observacion ingresada aparece al consultar el detalle de la tarea.
\item \textbf{Evidencia textual:} ``hay alguna informacion que considere que falta. Si, las observaciones de las tareas'' (EV-09).
\end{itemize}

\rf{RF-29}{Historial completo por parcela}
\begin{itemize}
\item \textbf{Descripcion:} el sistema debe ofrecer una vista de historial por parcela que agregue actividades, tratamientos y cosechas a lo largo del tiempo, no solo el estado actual.
\item \textbf{Actor y origen:} Tecnica. EV-16.
\item \textbf{Entradas:} identificador de parcela.
\item \textbf{Salidas:} linea de tiempo de eventos de la parcela.
\item \textbf{Precondicion:} la parcela tiene al menos un evento registrado.
\item \textbf{Postcondicion:} ninguna.
\item \textbf{Flujo principal:}
\begin{enumerate}[label=\arabic*.,leftmargin=1.6em]
\item El usuario selecciona una parcela y abre su vista de historial.
\item El sistema agrega actividades, tratamientos y cosechas de esa parcela a lo largo del tiempo.
\item El sistema muestra los eventos ordenados cronologicamente.
\end{enumerate}
\item \textbf{Flujos alternativos y excepciones:}
\begin{itemize}
\item Si la parcela no tiene ningun evento registrado, el sistema muestra el historial vacio en vez de un error (precondicion de RF-29).
\end{itemize}
\item \textbf{Prioridad MoSCoW:} Should have.
\item \textbf{Criterio de verificacion:} al consultar el historial de una parcela con 5 eventos registrados en distintas fechas, los 5 aparecen ordenados cronologicamente.
\item \textbf{Evidencia textual:} ``El historial de cada parcela'' (EV-16).
\end{itemize}

\rf{RF-30}{Umbral de stock minimo configurable por insumo}
\begin{itemize}
\item \textbf{Descripcion:} el sistema debe permitir definir, para cada insumo, un valor numerico de stock minimo que dispare la alerta de bajo stock del RF-08.
\item \textbf{Actor y origen:} Tecnica. EV-16.
\item \textbf{Entradas:} cantidad minima por insumo.
\item \textbf{Salidas:} umbral guardado en la ficha del insumo.
\item \textbf{Precondicion:} insumo ya registrado.
\item \textbf{Postcondicion:} el umbral queda disponible para el mecanismo de alertas.
\item \textbf{Flujo principal:}
\begin{enumerate}[label=\arabic*.,leftmargin=1.6em]
\item El Tecnico abre la ficha de un insumo ya registrado.
\item Define el valor numerico de stock minimo (stock limit) para ese insumo.
\item El sistema guarda el umbral y lo deja disponible para el mecanismo de alertas (RF-08).
\end{enumerate}
\item \textbf{Flujos alternativos y excepciones:}
\begin{itemize}
\item El Tecnico puede editar el umbral en cualquier momento; el nuevo valor se aplica desde la siguiente actualizacion de stock.
\end{itemize}
\item \textbf{Prioridad MoSCoW:} Should have.
\item \textbf{Criterio de verificacion:} al bajar el stock del insumo por debajo del valor configurado, se dispara la alerta del RF-08.
\item \textbf{Evidencia textual:} ``Sobre cada insumo, lo que es la cantidad, el nombre y ademas del stock limit'' (EV-16).
\end{itemize}

\rf{RF-31}{Deteccion de valores atipicos al ingresar datos}
\begin{itemize}
\item \textbf{Descripcion:} el sistema debe comparar cada dato numerico recien ingresado, por ejemplo la cantidad cosechada, contra el promedio historico del mismo lote y avisar si el valor se aleja de forma significativa, antes de guardarlo.
\item \textbf{Actor y origen:} Tecnico. EV-17.
\item \textbf{Entradas:} valor ingresado, historico del lote.
\item \textbf{Salidas:} advertencia de valor atipico.
\item \textbf{Precondicion:} existe historico suficiente del lote, minimo 3 registros previos.
\item \textbf{Postcondicion:} el usuario confirma o corrige el dato antes de guardar.
\item \textbf{Flujo principal:}
\begin{enumerate}[label=\arabic*.,leftmargin=1.6em]
\item El usuario ingresa un dato numerico (por ejemplo, cantidad cosechada) para un lote con historico suficiente (minimo 3 registros previos).
\item El sistema compara el valor contra el promedio historico del mismo lote.
\item Si el valor se desvia mas de 2 desviaciones estandar del promedio, el sistema muestra una advertencia antes de guardar.
\end{enumerate}
\item \textbf{Flujos alternativos y excepciones:}
\begin{itemize}
\item El usuario confirma el dato pese a la advertencia (si el valor atipico es real), o lo corrige antes de guardar (postcondicion de RF-31).
\item Si el lote no tiene historico suficiente, el sistema guarda el dato sin comparar (precondicion de RF-31).
\end{itemize}
\item \textbf{Prioridad MoSCoW:} Could have.
\item \textbf{Criterio de verificacion:} un valor que se desvia mas de 2 desviaciones estandar del promedio historico dispara la advertencia antes de guardar.
\item \textbf{Evidencia textual:} ``que al momento de meter un dato, una variable, me indique... Que este valor de aqui estuvo mal'' (EV-17).
\end{itemize}

\rf{RF-32}{Via de integracion futura con sensores de campo}
\begin{itemize}
\item \textbf{Descripcion:} el sistema debe exponer una interfaz capaz de recibir automaticamente lecturas de sensores de clima, suelo o nutrientes, como alternativa al ingreso manual, para una fase posterior del proyecto.
\item \textbf{Actor y origen:} Tecnico. EV-17.
\item \textbf{Entradas:} lectura de sensor, formato a definir en fase posterior.
\item \textbf{Salidas:} dato incorporado al historial del lote sin intervencion manual.
\item \textbf{Precondicion:} sensor compatible configurado.
\item \textbf{Postcondicion:} ninguna.
\item \textbf{Prioridad MoSCoW:} Won't have en esta version, queda documentado para trabajo futuro.
\item \textbf{Criterio de verificacion:} no aplica en esta version, queda como extension documentada.
\item \textbf{Evidencia textual:} ``Que ya la aplicacion que estan generando tenga sensores... Que tomen esos datos automaticamente'' (EV-17).
\end{itemize}

\rf{RF-33}{Bandeja de sugerencias para disputar una recomendacion de IA}
\begin{itemize}
\item \textbf{Descripcion:} el sistema debe ofrecer un mecanismo para que el usuario registre formalmente su desacuerdo con una recomendacion de IA, distinto del simple confirmar o descartar del RF-09, y que quede disponible como registro de retroalimentacion.
\item \textbf{Actor y origen:} Tecnico. EV-17.
\item \textbf{Entradas:} recomendacion en cuestion, comentario del usuario.
\item \textbf{Salidas:} registro de desacuerdo asociado a la recomendacion.
\item \textbf{Precondicion:} existe una recomendacion de IA previa.
\item \textbf{Postcondicion:} el desacuerdo queda visible para revision posterior del equipo tecnico.
\item \textbf{Flujo principal:}
\begin{enumerate}[label=\arabic*.,leftmargin=1.6em]
\item El usuario, en desacuerdo con una recomendacion de IA previa (RF-09/RF-10), abre la bandeja de sugerencias.
\item Registra su comentario asociado a esa recomendacion.
\item El sistema guarda el desacuerdo, visible para revision posterior del equipo tecnico.
\end{enumerate}
\item \textbf{Flujos alternativos y excepciones:}
\begin{itemize}
\item El desacuerdo registrado aparece en un listado consultable con fecha y comentario del usuario (criterio de verificacion de RF-33), sin alterar el estado de la alerta original.
\end{itemize}
\item \textbf{Prioridad MoSCoW:} Should have.
\item \textbf{Criterio de verificacion:} un desacuerdo registrado aparece en un listado consultable, con la fecha y el comentario del usuario.
\item \textbf{Evidencia textual:} ``yo obviamente que daria mi observacion en una bandejita de sugerencias'' (EV-17).
\end{itemize}

\rf{RF-34}{Esquema de atributos configurable por variedad de cultivo}
\begin{itemize}
\item \textbf{Descripcion:} el catalogo de cultivos del RF-02 debe permitir definir campos y variables distintas por variedad, no solo por tipo de cultivo, dado que el platano y las distintas variedades de cacao manejan datos de manejo agronomico diferentes.
\item \textbf{Actor y origen:} Tecnico. EV-17.
\item \textbf{Entradas:} variedad seleccionada.
\item \textbf{Salidas:} formulario con los campos correspondientes a esa variedad.
\item \textbf{Precondicion:} variedad dada de alta en el catalogo con su conjunto de atributos.
\item \textbf{Postcondicion:} ninguna.
\item \textbf{Flujo principal:}
\begin{enumerate}[label=\arabic*.,leftmargin=1.6em]
\item El Tecnico da de alta una variedad de cultivo en el catalogo (RF-02), definiendo el conjunto de campos y atributos propios de esa variedad.
\item Al seleccionar esa variedad en un formulario (por ejemplo, CU-01), el sistema muestra el conjunto de campos correspondiente.
\end{enumerate}
\item \textbf{Flujos alternativos y excepciones:}
\begin{itemize}
\item Al seleccionar una variedad de cacao, el formulario muestra un conjunto de campos distinto al que se muestra al seleccionar platano (criterio de verificacion de RF-34).
\end{itemize}
\item \textbf{Prioridad MoSCoW:} Must have.
\item \textbf{Criterio de verificacion:} al seleccionar una variedad de cacao, el formulario muestra un conjunto de campos distinto al que se muestra al seleccionar platano.
\item \textbf{Evidencia textual:} ``Existen muchas diferencias, en el platano se registra lo que es... En el cacao se registra desde su mazorca... Hay muchos tipos de cacao, el nacional, el CN51'' (EV-17).
\end{itemize}

\rf{RF-35}{Registro de analisis de suelo previo a la siembra}
\begin{itemize}
\item \textbf{Descripcion:} el sistema debe permitir registrar, por parcela, el resultado de un analisis de suelo previo al establecimiento de un cultivo, indicando tipo de suelo y aptitud.
\item \textbf{Criterio legal cubierto:} C21 (Res. AGROCALIDAD 0072, Art. 8 Res.183).
\item \textbf{Actor y origen:} Tecnico. EV-17.
\item \textbf{Entradas:} tipo de suelo, resultado de aptitud, fecha del analisis.
\item \textbf{Salidas:} ficha de analisis de suelo asociada a la parcela.
\item \textbf{Precondicion:} parcela registrada.
\item \textbf{Postcondicion:} dato consultable antes de aprobar la siembra.
\item \textbf{Flujo principal:}
\begin{enumerate}[label=\arabic*.,leftmargin=1.6em]
\item El Tecnico registra, para una parcela ya existente, el resultado de un analisis de suelo (calicata) previo a la siembra: tipo de suelo, aptitud y fecha.
\item El sistema guarda la ficha de analisis de suelo asociada a la parcela.
\end{enumerate}
\item \textbf{Flujos alternativos y excepciones:}
\begin{itemize}
\item Una parcela con analisis de suelo registrado muestra el resultado de aptitud al consultar su ficha, antes de aprobar la siembra (criterio de verificacion de RF-35).
\end{itemize}
\item \textbf{Prioridad MoSCoW:} Could have.
\item \textbf{Criterio de verificacion:} una parcela con analisis de suelo registrado muestra el resultado de aptitud al consultar su ficha.
\item \textbf{Evidencia textual:} ``una calicata, para ver o realizar un analisis de suelo... nos va a decir si es apto o no es apto sembrar en dicho terreno'' (EV-17).
\end{itemize}

\subsection{Requisitos legales derivados}

Cada requisito derivado de norma lleva en su ficha el campo \textbf{Criterio legal cubierto}, con el identificador propio del criterio (C1 a C26), la norma de la que procede y el articulo exacto. Los 26 criterios son los definidos en \ruta{01\_ERS/Modelo\_Legal\_LOPDP.md} y desarrollados en \ruta{01\_ERS/Matriz\_Obligacion\_Legal.md}, y son los mismos que evalua el componente empirico en \ruta{07\_Datos/datos\_crudos/cobertura\_legal.csv}. De este modo el identificador legal no vive solo en la matriz de trazabilidad: viaja con el requisito dentro de este documento, y cualquier lector puede ir del requisito al articulo sin salir del ERS.


Estos requisitos se derivan directamente de los 26 criterios, C1 a C26, de \ruta{Modelo\_Legal\_LOPDP.md}, en sus tres bloques normativos. Ninguna de las 17 entrevistas menciono el proceso de certificacion BPA, el derecho de acceso o rectificacion de datos personales, el aviso fitosanitario formal ni el registro de capacitaciones, asi que siguen siendo vacios legales puros. El equipo de proteccion personal y el analisis de suelo si encontraron respaldo parcial en EV-17, ver RF-18 y RF-35, pero el certificado de salud del trabajador continua sin evidencia de entrevista.

\rf{RF-22}{Registro de consentimiento del trabajador para el tratamiento de sus datos personales}
\begin{itemize}
\item \textbf{Actor y origen:} Administrador. Sin evidencia de entrevista; requisito derivado de norma (ver criterio legal cubierto).
\item \textbf{Criterio legal cubierto:} C1 (LOPDP, Art. 4, 12(8)); C2 (LOPDP, Art. 4, 34-35); C3 (LOPDP, Art. 10(d)); C5 (LOPDP, Art. 12(5)); C12 (LOPDP, Art. 34-35).
\item \textbf{Flujo principal:}
\begin{enumerate}[label=\arabic*.,leftmargin=1.6em]
\item En el primer inicio de sesion de un trabajador (CU-09), el sistema muestra el aviso de tratamiento de datos personales (LOPDP Art. 8).
\item El usuario acepta el consentimiento de forma libre, especifica, informada e inequivoca.
\item El sistema registra el consentimiento antes de continuar con cualquier otro modulo.
\end{enumerate}
\item \textbf{Flujos alternativos y excepciones:}
\begin{itemize}
\item Si el usuario no acepta el aviso, el sistema no concede acceso a ningun modulo que trate datos personales (regla de negocio de CU-09).
\end{itemize}
\item \textbf{Prioridad MoSCoW:} Must have, regulatorio.
\item \textbf{Evidencia legal:} Art. 8 LOPDP. El consentimiento debe ser libre, especifico, informado e inequivoco.
\end{itemize}

\rf{RF-23}{Modulo de derechos ARCO+ del trabajador}
\begin{itemize}
\item \textbf{Descripcion:} acceso, rectificacion y eliminacion de sus propios datos.
\item \textbf{Criterio legal cubierto:} C4 (LOPDP, Art. 8); C7 (LOPDP, Art. 10(i)); C8 (LOPDP, Art. 33, 36); C9 (LOPDP, Art. 43, 46); C10 (LOPDP, Art. 13-19).
\item \textbf{Actor y origen:} Jornalero o trabajador. Sin evidencia de entrevista, criterio C10, LOPDP Art. 13 a 19.
\item \textbf{Flujo principal:}
\begin{enumerate}[label=\arabic*.,leftmargin=1.6em]
\item El trabajador, autenticado en el sistema, solicita acceder, rectificar o eliminar sus propios datos personales desde Ajustes → Mis Datos.
\item El sistema atiende la solicitud sobre los datos del propio usuario.
\item El sistema deja un registro de la accion realizada (criterio de verificacion derivado de C10, LOPDP Art. 13-19).
\end{enumerate}
\item \textbf{Flujos alternativos y excepciones:}
\begin{itemize}
\item Si la solicitud es de eliminacion y existen datos que deben conservarse por obligacion legal (por ejemplo, trazabilidad BPA), el sistema informa al usuario cuales datos no pueden eliminarse y por que.
\end{itemize}
\item \textbf{Prioridad MoSCoW:} Must have, regulatorio.
\end{itemize}

\rf{RF-24}{Registro de certificado de salud del trabajador que aplica agroquimicos}
\begin{itemize}
\item \textbf{Actor y origen:} Administrador. Sin evidencia de entrevista para el certificado de salud, criterio C17, Resolucion AGROCALIDAD 183, Art. 33 y 34. El requisito de equipo de proteccion personal ya cuenta con evidencia real en RF-18.
\item \textbf{Criterio legal cubierto:} C17 (Res. AGROCALIDAD 183, Art. 31-34).
\item \textbf{Flujo principal:}
\begin{enumerate}[label=\arabic*.,leftmargin=1.6em]
\item Antes de asignar una tarea de aplicacion de agroquimicos (CU-11), el Administrador registra el certificado de salud vigente del trabajador.
\item El sistema adjunta el certificado al perfil del trabajador.
\end{enumerate}
\item \textbf{Flujos alternativos y excepciones:}
\begin{itemize}
\item Si el trabajador no tiene certificado de salud vigente registrado, el sistema advierte que el certificado no esta registrado al intentar asignarle esa tarea (Res. AGROCALIDAD 183, Art. 33-34).
\end{itemize}
\item \textbf{Prioridad MoSCoW:} Must have, regulatorio.
\end{itemize}

\rf{RF-25}{Registro y seguimiento del proceso de certificacion BPA ante AGROCALIDAD}
\begin{itemize}
\item \textbf{Descripcion:} solicitud, inspeccion y vigencia.
\item \textbf{Criterio legal cubierto:} C20 (Res. AGROCALIDAD 183, Art. 39-43).
\item \textbf{Actor y origen:} Administrador. Sin evidencia de entrevista, criterio C20, Resolucion AGROCALIDAD 183, Art. 39 a 43.
\item \textbf{Flujo principal:}
\begin{enumerate}[label=\arabic*.,leftmargin=1.6em]
\item El Administrador revisa el panel de cumplimiento (vigencia del certificado BPA, capacitaciones pendientes, bitacora de bioseguridad) (CU-13).
\item Si el certificado esta por vencer o falta evidencia de respaldo, solicita la renovacion ante AGROCALIDAD.
\item El sistema adjunta como evidencia los registros de capacitacion y la bitacora de bioseguridad.
\end{enumerate}
\item \textbf{Flujos alternativos y excepciones:}
\begin{itemize}
\item Si falta una capacitacion obligatoria, el sistema bloquea la solicitud de renovacion hasta que se registre (RF-36, RF-39).
\end{itemize}
\item \textbf{Prioridad MoSCoW:} Should have, regulatorio.
\end{itemize}

\rf{RF-36}{Registro de capacitacion en manejo de plaguicidas y primeros auxilios}
\begin{itemize}
\item \textbf{Actor y origen:} Administrador. Sin evidencia de entrevista, criterio C23, Resolucion AGROCALIDAD 183, Art. 18(e).
\item \textbf{Criterio legal cubierto:} C23 (Res. AGROCALIDAD 0072, Art. 18(e) Res.183).
\item \textbf{Flujo principal:}
\begin{enumerate}[label=\arabic*.,leftmargin=1.6em]
\item El Administrador registra una capacitacion en manejo de plaguicidas o primeros auxilios recibida por un trabajador.
\item El sistema guarda el registro como evidencia de respaldo para la certificacion BPA (RF-25, CU-13).
\end{enumerate}
\item \textbf{Flujos alternativos y excepciones:}
\begin{itemize}
\item Si la capacitacion registrada esta vencida o incompleta, el sistema la marca como pendiente en el panel de cumplimiento (CU-13).
\end{itemize}
\item \textbf{Prioridad MoSCoW:} Should have, regulatorio.
\end{itemize}

\rf{RF-37}{Aviso ante sintomas sospechosos de plaga cuarentenaria}
\begin{itemize}
\item \textbf{Actor y origen:} Jornalero, Tecnico, Administrador. Sin evidencia de entrevista, criterio C24, Resolucion AGROCALIDAD 0072, Art. 3.6.1(a). EV-17 menciono el problema del Moko y las medidas de desinfeccion, pero ningun entrevistado describio el mecanismo formal de aviso a AGROCALIDAD.
\item \textbf{Criterio legal cubierto:} C24 (Res. AGROCALIDAD 0072, Art. 3.6.1(a) Res.0072).
\item \textbf{Flujo principal:}
\begin{enumerate}[label=\arabic*.,leftmargin=1.6em]
\item Un usuario de campo detecta sintomas sospechosos de plaga cuarentenaria (ej. Moko) en un lote.
\item Registra el sintoma y el lote afectado en el sistema.
\item El sistema genera el aviso fitosanitario y lo marca con prioridad alta, visible de inmediato para el Administrador.
\item Un Tecnico confirma el sintoma tras la inspeccion.
\end{enumerate}
\item \textbf{Flujos alternativos y excepciones:}
\begin{itemize}
\item El Tecnico descarta el sintoma tras la inspeccion; el aviso no se envia a AGROCALIDAD, pero queda registrado como descartado para trazabilidad (CU-14).
\end{itemize}
\item \textbf{Prioridad MoSCoW:} Must have, regulatorio.
\end{itemize}

\rf{RF-38}{Bitacora de bioseguridad de ingreso y salida del predio}
\begin{itemize}
\item \textbf{Actor y origen:} Administrador. Sin evidencia de entrevista sobre el registro formal, aunque EV-17 describio la practica de desinfeccion. Criterio C25, Resolucion AGROCALIDAD 0072, Art. 3.6.1(d).
\item \textbf{Criterio legal cubierto:} C25 (Res. AGROCALIDAD 0072, Art. 3.6.1(d) Res.0072).
\item \textbf{Flujo principal:}
\begin{enumerate}[label=\arabic*.,leftmargin=1.6em]
\item Al confirmarse una visita al predio (tecnica o por sintoma sospechoso, CU-10/CU-14), el usuario registra el evento de ingreso/salida en la bitacora de bioseguridad.
\item El sistema guarda el registro con fecha y motivo de la visita.
\end{enumerate}
\item \textbf{Flujos alternativos y excepciones:}
\begin{itemize}
\item La bitacora queda disponible como evidencia de respaldo para la certificacion BPA y para la auditoria de AGROCALIDAD (RF-25, CU-13).
\end{itemize}
\item \textbf{Prioridad MoSCoW:} Should have, regulatorio.
\end{itemize}

\rf{RF-39}{Registro de capacitaciones del personal sobre control fitosanitario especifico}
\begin{itemize}
\item \textbf{Actor y origen:} Administrador. Sin evidencia de entrevista, criterio C26, Resolucion AGROCALIDAD 0072, Art. 3.6.1(f)(g).
\item \textbf{Criterio legal cubierto:} C26 (Res. AGROCALIDAD 0072, Art. 3.6.1(f)(g) Res.0072).
\item \textbf{Flujo principal:}
\begin{enumerate}[label=\arabic*.,leftmargin=1.6em]
\item El Administrador registra una capacitacion fitosanitaria especifica recibida por un trabajador (control de plagas cuarentenarias, bioseguridad).
\item El sistema guarda el registro y lo deja disponible como evidencia de auditoria (RF-25, CU-13).
\end{enumerate}
\item \textbf{Flujos alternativos y excepciones:}
\begin{itemize}
\item Si falta una capacitacion fitosanitaria obligatoria, el sistema bloquea la solicitud de renovacion BPA hasta que se registre (regla de negocio de CU-13).
\end{itemize}
\item \textbf{Prioridad MoSCoW:} Should have, regulatorio.
\end{itemize}

\subsection{Requisitos no funcionales}

\rf{RNF-01}{Acceso movil prioritario}
\begin{itemize}
\item \textbf{Categoria ISO/IEC 25010:} interaccion con el usuario.
\item \textbf{Descripcion:} las funciones de consulta y registro mas usadas en campo, tareas, cosecha e inventario, deben estar disponibles desde un telefono movil, no solo desde escritorio.
\item \textbf{Metrica:} el 100\% de las pantallas de RF-01, RF-03, RF-04, RF-11 y RF-19 deben renderizar correctamente en una pantalla de 360 por 800 pixeles sin desplazamiento horizontal.
\item \textbf{Evidencia textual:} ``Desde mi telefono, mas rapido'' (EV-01).
\end{itemize}

\rf{RNF-02}{Interfaz simple y sin tecnicismos}
\begin{itemize}
\item \textbf{Categoria ISO/IEC 25010:} interaccion con el usuario.
\item \textbf{Descripcion:} las etiquetas de campos y menus deben usar vocabulario cotidiano del agricultor, evitando abreviaturas o nombres de columna tecnicos.
\item \textbf{Metrica:} en una prueba de usabilidad con 5 jornaleros sin experiencia previa en el sistema, al menos el 90\% de los terminos de la interfaz deben ser comprendidos sin ayuda externa.
\item \textbf{Evidencia textual:} ``cosas muy tecnicas, numeros de oficina, codigos raros a la vez, solo confundir. Entre mas directo y sencillo es mejor'' (EV-15).
\end{itemize}

\rf{RNF-03}{Usabilidad para usuarios con baja alfabetizacion digital}
\begin{itemize}
\item \textbf{Categoria ISO/IEC 25010:} interaccion con el usuario.
\item \textbf{Descripcion:} un usuario sin experiencia previa en aplicaciones moviles debe poder completar el registro de una cosecha sin asistencia despues de una capacitacion breve.
\item \textbf{Metrica:} tiempo maximo de 5 minutos de capacitacion para que un jornalero sin experiencia previa registre una cosecha sin errores.
\item \textbf{Evidencia textual:} ``No tengo estudios'' y uso previo de aplicaciones ``No'' (EV-02).
\end{itemize}

\rf{RNF-04}{Robustez de campo ante condiciones adversas}
\begin{itemize}
\item \textbf{Categoria ISO/IEC 25010:} fiabilidad.
\item \textbf{Descripcion:} el registro de actividades y cosecha debe seguir siendo posible sin conexion a internet, sincronizando los datos cuando la conexion se restablezca.
\item \textbf{Metrica:} el sistema debe permitir guardar localmente hasta 7 dias de registros sin conexion, sin perdida de datos al reconectarse.
\item \textbf{Evidencia textual:} ``el barro y la lluvia... Danan las hojas de papel donde nosotros anotamos y nos confundimos'' (EV-14).
\end{itemize}

\rf{RNF-05}{Explicabilidad verificable de las recomendaciones de IA}
\begin{itemize}
\item \textbf{Categoria ISO/IEC 25010:} adecuacion funcional.
\item \textbf{Descripcion:} toda recomendacion de IA debe mostrar la razon de la sugerencia antes de que el usuario pueda aplicarla.
\item \textbf{Metrica:} el 100\% de las alertas de IA, RF-09 y RF-10, deben incluir un texto de justificacion de un maximo de 60 palabras, visible antes del boton de confirmacion.
\item \textbf{Evidencia textual:} ``Que explique cual es el problema y lo que se debe hacer'' (EV-09).
\end{itemize}

\rf{RNF-06}{Transparencia del origen de los datos de la recomendacion}
\begin{itemize}
\item \textbf{Categoria ISO/IEC 25010:} adecuacion funcional.
\item \textbf{Descripcion:} ademas de explicar el motivo, la recomendacion debe indicar en que datos del propio lote se baso.
\item \textbf{Metrica:} el 100\% de las recomendaciones deben listar al menos 1 dato de origen, por ejemplo el registro de fumigacion del 12 de agosto.
\item \textbf{Evidencia textual:} ``Los datos en que se basa esta informacion'' (EV-16).
\end{itemize}

\rf{RNF-07}{Preservacion del historico de datos sin perdida}
\begin{itemize}
\item \textbf{Categoria ISO/IEC 25010:} fiabilidad.
\item \textbf{Descripcion:} ningun dato de cosecha, actividad o inventario registrado debe poder eliminarse de forma permanente sin confirmacion explicita y registro de auditoria.
\item \textbf{Metrica:} el sistema debe mantener un registro de auditoria de cambios y eliminaciones por al menos 24 meses.
\item \textbf{Evidencia textual:} ``habia desvarianza de datos, en lo que no se podia cuadrar un numero de cajas'' (EV-01).
\end{itemize}

\rf{RNF-08}{Consentimiento explicito registrado en el primer inicio de sesion}
\begin{itemize}
\item \textbf{Categoria ISO/IEC 25010:} seguridad.
\item \textbf{Criterio legal cubierto:} C6 (LOPDP, Art. 37, 41); C13 (LOPDP, Art. 47(2)(3)(14)).
\item \textbf{Descripcion:} todo trabajador debe aceptar un aviso de tratamiento de datos personales antes de poder usar el sistema por primera vez.
\item \textbf{Metrica:} el 100\% de las cuentas activas deben tener un registro de aceptacion con fecha y hora.
\item \textbf{Evidencia legal:} Art. 8 LOPDP.
\end{itemize}

\rf{RNF-09}{Seguridad de los datos personales almacenados}
\begin{itemize}
\item \textbf{Categoria ISO/IEC 25010:} seguridad, derivado legal.
\item \textbf{Criterio legal cubierto:} C6 (LOPDP, Art. 37, 41).
\item \textbf{Descripcion:} los datos personales de los trabajadores deben cifrarse en reposo y en transito.
\item \textbf{Metrica:} cifrado AES-256 en base de datos y TLS 1.2 o superior en toda comunicacion cliente y servidor.
\item \textbf{Evidencia legal:} Art. 37 y 41 LOPDP, sin evidencia directa de entrevista, requisito regulatorio.
\end{itemize}

\rf{RNF-10}{Posibilidad de disentir de una recomendacion sin bloquear el flujo de trabajo}
\begin{itemize}
\item \textbf{Categoria ISO/IEC 25010:} interaccion con el usuario.
\item \textbf{Descripcion:} el usuario debe poder continuar su trabajo normalmente aunque decida no seguir una recomendacion de IA, sin que el sistema se lo impida ni lo penalice.
\item \textbf{Metrica:} cero casos en los que el sistema bloquee una accion del usuario por no aceptar una recomendacion de IA.
\item \textbf{Evidencia textual:} ``Lo primero que en este caso haria es comparar con mi experiencia'' (EV-16).
\end{itemize}

\rf{RNF-11}{Reduccion del esfuerzo de captura manual de datos}
\begin{itemize}
\item \textbf{Categoria ISO/IEC 25010:} eficiencia de desempeno.
\item \textbf{Descripcion:} las pantallas de registro frecuente, cosecha y actividad, deben minimizar el numero de campos obligatorios y usar seleccion en vez de texto libre donde sea posible.
\item \textbf{Metrica:} registrar una cosecha completa debe tomar un maximo de 30 segundos en condiciones normales de campo.
\item \textbf{Evidencia textual:} ``los datos de manera manual... eso si se hace un poco tedioso'' (EV-17).
\end{itemize}

\rf{RNF-12}{Conservacion minima de registros de trazabilidad}
\begin{itemize}
\item \textbf{Categoria ISO/IEC 25010:} fiabilidad, derivado legal.
\item \textbf{Descripcion:} los registros de actividad, cosecha y aplicacion de insumos deben conservarse un minimo de 2 anos, conforme a la normativa de trazabilidad agroindustrial.
\item \textbf{Metrica:} ningun registro de los tipos mencionados puede eliminarse antes de cumplir 24 meses desde su creacion.
\item \textbf{Evidencia legal:} Art. 38 de la Resolucion AGROCALIDAD 183, sin evidencia directa de entrevista, requisito regulatorio.
\end{itemize}

\rf{RNF-13}{Compatibilidad con dispositivos de gama baja}
\begin{itemize}
\item \textbf{Categoria ISO/IEC 25010:} compatibilidad.
\item \textbf{Descripcion:} la aplicacion movil debe funcionar en equipos Android de gama baja, dado el perfil economico de los usuarios de campo.
\item \textbf{Metrica:} funcionamiento fluido, sin caidas, en un dispositivo con 2 GB de RAM y Android 9 o superior.
\item \textbf{Evidencia textual:} inferido de que ningun jornalero entrevistado menciono tener un telefono de alta gama. Varios reportaron bajo uso previo de aplicaciones, EV-02 y EV-13.
\end{itemize}

\rf{RNF-14}{Extensibilidad del catalogo de cultivos y variedades}
\begin{itemize}
\item \textbf{Categoria ISO/IEC 25010:} mantenibilidad.
\item \textbf{Descripcion:} agregar un nuevo cultivo o variedad al catalogo, RF-02 y RF-34, no debe requerir cambios en el codigo de la aplicacion, solo configuracion.
\item \textbf{Metrica:} un administrador debe poder agregar una variedad nueva con su conjunto de atributos en menos de 10 minutos, sin intervencion del equipo de desarrollo.
\item \textbf{Evidencia textual:} ``seria bueno que lo inserten... Maracuya, la pimienta y el cafe'' (EV-10).
\end{itemize}

\rf{RNF-15}{Disponibilidad del servicio de notificaciones}
\begin{itemize}
\item \textbf{Categoria ISO/IEC 25010:} fiabilidad.
\item \textbf{Descripcion:} las notificaciones de bajo stock y asignacion de tareas deben entregarse de forma confiable incluso con conectividad intermitente.
\item \textbf{Metrica:} el 95\% de las notificaciones deben entregarse dentro de los 5 minutos posteriores a la reconexion del dispositivo.
\item \textbf{Evidencia textual:} ``Mediante una notificacion como un mensaje de texto, algo que creo que es lo mas facil'' (EV-10).
\end{itemize}

\subsubsection{Requisitos no funcionales del componente de inteligencia artificial}

Los RF-09, RF-10, RF-31 y RF-33 introducen un componente de IA (alerta de plagas/enfermedades y diagnostico por imagen). Las seis RNF siguientes lo gobiernan de forma transversal, con el formato identificador, metrica, unidad, umbral, metodo de verificacion, responsable y frecuencia acordado en el reparto de roles del equipo del 2026-09-02.

\rf{RNF-16}{Desempeno de deteccion/prediccion del modelo de IA}
\begin{itemize}
\item \textbf{Categoria ISO/IEC 25010:} fiabilidad del componente de IA.
\item \textbf{Descripcion:} el modelo que sustenta RF-09 y RF-10 debe alcanzar un nivel minimo de acierto en la deteccion de plagas/enfermedades antes de habilitarse en produccion.
\item \textbf{Metrica:} exactitud (accuracy) y sensibilidad (recall) del modelo sobre el conjunto de validacion.
\item \textbf{Unidad:} porcentaje (\%).
\item \textbf{Umbral:} $\geq$80\% de exactitud y $\geq$75\% de sensibilidad en el conjunto de validacion, antes de habilitar el modelo en campo.
\item \textbf{Metodo de verificacion:} evaluacion del modelo sobre un conjunto de prueba etiquetado, ejecutada por script versionado, con matriz de confusion documentada en el repositorio.
\item \textbf{Responsable:} equipo tecnico (especificacion: Danela Arteaga; analisis estadistico: Maria Escudero).
\item \textbf{Donde interviene:} CU-06, paso 1 del flujo principal (donde el modelo genera la sugerencia); tambien en la excepcion CU-06.1e cuando se usa diagnostico por imagen (RF-10).
\end{itemize}

\rf{RNF-17}{Explicabilidad de las recomendaciones de IA}
\begin{itemize}
\item \textbf{Categoria ISO/IEC 25010:} transparencia.
\item \textbf{Descripcion:} toda alerta o recomendacion generada por el componente de IA (RF-09, RF-10) debe mostrar al usuario al menos un factor que motivo la sugerencia, no solo el resultado, para que decida con criterio si confirmarla (RF-09) o disputarla (RF-33).
\item \textbf{Metrica:} proporcion de alertas emitidas que muestran al menos un factor explicativo.
\item \textbf{Unidad:} porcentaje (\%) de alertas con explicacion.
\item \textbf{Umbral:} 100\% de las alertas deben mostrar al menos un factor explicativo.
\item \textbf{Metodo de verificacion:} revision de una muestra de alertas generadas en pruebas de aceptacion, confirmando la presencia del campo de explicacion.
\item \textbf{Responsable:} Danela Arteaga (especificacion); equipo tecnico (implementacion).
\item \textbf{Donde interviene:} CU-06, pasos 1 y 2 del flujo principal (donde se genera la justificacion y se notifica al usuario marcada como sugerencia).
\end{itemize}

\rf{RNF-18}{Equidad de desempeno entre cultivos}
\begin{itemize}
\item \textbf{Categoria ISO/IEC 25010:} equidad.
\item \textbf{Descripcion:} el componente de IA no debe generar sistematicamente peor desempeno (mas falsos negativos) para un cultivo respecto del otro, dado que cacao y platano estan igualmente en alcance del estudio de caso.
\item \textbf{Metrica:} diferencia en la tasa de falsos negativos entre cacao y platano.
\item \textbf{Unidad:} puntos porcentuales de diferencia.
\item \textbf{Umbral:} diferencia $\leq$10 puntos porcentuales en la tasa de falsos negativos entre ambos cultivos.
\item \textbf{Metodo de verificacion:} evaluacion separada del modelo por subconjunto (cacao, platano) sobre el conjunto de validacion, con reporte comparativo versionado.
\item \textbf{Responsable:} equipo tecnico; revision independiente por Maria Escudero (analisis estadistico).
\item \textbf{Donde interviene:} CU-06, paso 1 del flujo principal (la deteccion del modelo, medida por separado para cacao y platano).
\end{itemize}

\rf{RNF-19}{Supervision humana obligatoria antes de cualquier accion automatica}
\begin{itemize}
\item \textbf{Categoria ISO/IEC 25010:} supervision humana.
\item \textbf{Descripcion:} formaliza como requisito transversal lo que RF-09 ya exige en su flujo: ninguna alerta o recomendacion de IA puede ejecutar una accion sobre los datos del sistema sin confirmacion explicita de una persona.
\item \textbf{Metrica:} numero de acciones aplicadas automaticamente sobre datos del sistema sin confirmacion humana.
\item \textbf{Unidad:} conteo de incidentes.
\item \textbf{Umbral:} 0 acciones automaticas sin confirmacion humana.
\item \textbf{Metodo de verificacion:} revision de logs de auditoria y prueba funcional dirigida a intentar forzar una aplicacion automatica.
\item \textbf{Responsable:} equipo tecnico; verificacion independiente por Kamila Calle (gatekeeper P11).
\item \textbf{Donde interviene:} CU-06, pasos 3 y 4 del flujo principal (revision y confirmacion humana); tambien en la excepcion CU-06.3e, donde la alerta queda fuera del conjunto aplicado hasta revision del equipo tecnico (RF-33).
\end{itemize}

\rf{RNF-20}{Monitoreo continuo del desempeno del modelo en produccion}
\begin{itemize}
\item \textbf{Categoria ISO/IEC 25010:} monitoreo.
\item \textbf{Descripcion:} el sistema debe registrar, para cada alerta de IA, si el usuario la confirmo, la descarto (RF-09) o la disputo (RF-33), para poder calcular el desempeno real del modelo en campo y no solo en el conjunto de prueba.
\item \textbf{Metrica:} proporcion de alertas emitidas con retroalimentacion registrada (confirmada, descartada o disputada).
\item \textbf{Unidad:} porcentaje (\%).
\item \textbf{Umbral:} $\geq$90\% de las alertas deben tener retroalimentacion registrada dentro de los 7 dias de emitidas.
\item \textbf{Metodo de verificacion:} consulta sobre el registro de alertas y su estado, generada por script versionado.
\item \textbf{Responsable:} Roselyn Sanchez (infraestructura de datos); Danela Arteaga (especificacion).
\item \textbf{Donde interviene:} CU-06, paso 4 del flujo principal (registro de la decision); tambien en el flujo alternativo CU-06.3a (descarte) y en la excepcion CU-06.3e (desacuerdo formal, RF-33).
\end{itemize}

\rf{RNF-21}{Clasificacion de riesgo de las recomendaciones de IA}
\begin{itemize}
\item \textbf{Categoria ISO/IEC 25010:} gestion de riesgo.
\item \textbf{Descripcion:} cada tipo de alerta o recomendacion del componente de IA debe clasificarse por nivel de riesgo (bajo, medio, alto) segun el impacto potencial de seguirla sin verificacion, y ese nivel debe ser visible al usuario junto a la alerta.
\item \textbf{Metrica:} proporcion de tipos de alerta definidos en RF-09/RF-10 con nivel de riesgo asignado y visible en la interfaz.
\item \textbf{Unidad:} porcentaje (\%) de tipos de alerta clasificados.
\item \textbf{Umbral:} 100\% de los tipos de alerta deben tener un nivel de riesgo asignado antes de habilitarse.
\item \textbf{Metodo de verificacion:} revision de la tabla de configuracion de tipos de alerta, confirmando el campo de nivel de riesgo.
\item \textbf{Responsable:} Danela Arteaga (especificacion); equipo tecnico (implementacion).
\item \textbf{Donde interviene:} CU-06, paso 2 del flujo principal (donde se notifica la alerta al usuario, junto con su nivel de riesgo visible).
\end{itemize}


\subsubsection{Clasificacion de riesgo del componente de IA y base legal}

La clasificacion de riesgo del componente de inteligencia artificial, la base legal que la fundamenta y el plan de monitoreo en produccion se desarrollan en el documento \ruta{01\_ERS/Clasificacion\_Riesgo\_IA.md}. Ese documento clasifica el sistema en dos niveles: como tratamiento de datos conserva la Categoria C, riesgo minimo operativo, porque el componente opera sobre material vegetal y datos operativos de la finca y no sobre datos personales ni biometricos; y por tipo de recomendacion asigna un nivel de riesgo a cada salida del componente, que es lo que exige RNF-21. Tres tipos quedan clasificados como riesgo alto (la sugerencia de aplicacion fitosanitaria de RF-09, la sospecha de plaga cuarentenaria de RF-09 y RF-37, y el diagnostico por imagen de RF-10) porque seguirlos sin verificacion puede causar dano economico o fitosanitario dificilmente reversible o activar una obligacion regulatoria ante AGROCALIDAD. La base legal directa de la explicabilidad exigida en RNF-17 es el Art. 20 de la LOPDP, que reconoce el derecho a no ser objeto de una decision basada unica o parcialmente en valoraciones automatizadas y a exigir una explicacion motivada. El plan de monitoreo desarrolla siete indicadores con umbral, frecuencia, responsable y accion correctiva. El componente esta especificado y no desplegado: no existe todavia un modelo entrenado, de modo que los umbrales de RNF-16 y RNF-18 operan como condiciones de habilitacion previas al despliegue, no como resultados obtenidos.

\subsection{Historias de usuario y criterios de aceptacion}

Los 17 requisitos funcionales de prioridad Must have se tradujeron a historias de usuario en formato Connextra, seguidas de un escenario de aceptacion en Gherkin construido a partir del criterio de verificacion ya definido para cada requisito. Cada historia cumple los criterios INVEST al describir una necesidad independiente, negociable en su implementacion, de valor claro para el usuario, estimable, acotada a un solo requisito y verificable mediante el escenario que la acompana.

\newcommand{\hu}[3]{\par\smallskip\noindent\textbf{#1, ligada a #2.} #3\par}

\hu{HU-01}{RF-01}{Como administrador, tecnico o jornalero, quiero registrar y editar una parcela con sus datos basicos, para tener siempre a la mano el nombre, el cultivo y la cantidad de plantas de cada lote sin depender de la memoria o de papeles sueltos.\\
\textit{Dado} que estoy autenticado con permiso de edicion, \textit{cuando} registro una parcela completando los siete campos obligatorios, \textit{entonces} el sistema la muestra en el listado general en menos de dos segundos y sin campos vacios.}

\hu{HU-02}{RF-03}{Como jornalero o tecnico, quiero registrar cada actividad agricola con su fecha, lote, producto y responsable, para dejar un historial confiable de lo que se hizo en cada parcela.\\
\textit{Dado} que estoy autenticado, \textit{cuando} intento guardar una actividad sin indicar la fecha, el lote o el trabajador responsable, \textit{entonces} el sistema rechaza el registro hasta que esos tres datos queden completos.}

\hu{HU-03}{RF-04}{Como jornalero, quiero registrar la cosecha en la unidad propia de cada cultivo, para no forzar el dato a una unidad que no corresponde con la forma real de trabajo en campo.\\
\textit{Dado} que estoy registrando la cosecha de un lote de cacao, \textit{cuando} elijo la unidad de medida, \textit{entonces} el sistema solo acepta tacho o quintal, y si el lote es de platano solo acepta racimo o caja.}

\hu{HU-04}{RF-05}{Como administrador, tecnico o jornalero, quiero ver el rendimiento promedio de cada lote por periodo, para comparar el desempeno de mis parcelas sin calcularlo a mano.\\
\textit{Dado} un lote con dos o mas registros de cosecha, \textit{cuando} consulto su rendimiento, \textit{entonces} el sistema muestra un promedio igual al que resulta de dividir la suma de las cosechas entre el numero de registros.}

\hu{HU-05}{RF-06}{Como tecnico o jornalero, quiero que el sistema calcule automaticamente la ganancia neta de un periodo, para conocer el resultado economico real sin hacer la resta manualmente.\\
\textit{Dado} que existen registros de ingresos y egresos en un periodo, \textit{cuando} consulto la ganancia neta, \textit{entonces} el valor mostrado es igual a la suma de ingresos menos la suma de egresos de ese periodo.}

\hu{HU-06}{RF-07}{Como administrador, tecnico o jornalero, quiero registrar la cantidad usada de un insumo, para saber en todo momento cuanto stock me queda sin contarlo manualmente.\\
\textit{Dado} un insumo con treinta unidades disponibles, \textit{cuando} registro un uso de veintiseis unidades, \textit{entonces} el sistema muestra cuatro unidades disponibles.}

\hu{HU-07}{RF-08}{Como administrador, jornalero o tecnica, quiero recibir una alerta cuando el stock de un insumo baje del umbral configurado, para comprar a tiempo antes de quedarme sin insumo para una labor programada.\\
\textit{Dado} un insumo cuyo stock cae por debajo del umbral configurado, \textit{cuando} se registra ese descenso, \textit{entonces} el sistema envia la notificacion al responsable en menos de cinco minutos.}

\hu{HU-08}{RF-09}{Como administrador, tecnico o jornalero, quiero recibir una alerta de posible plaga o enfermedad que siempre pida mi confirmacion antes de aplicarse, para mantener el control sobre las decisiones de manejo del cultivo.\\
\textit{Dado} que el sistema genera una alerta de plaga sobre un lote, \textit{cuando} la alerta se muestra al usuario responsable, \textit{entonces} queda registrada como pendiente de confirmar o descartar y en ningun caso se aplica de forma automatica al lote.}

\hu{HU-09}{RF-11}{Como administrador, tecnico o jornalero, quiero asignar tareas a un trabajador por lote y fecha y poder filtrarlas por estado, para saber en cualquier momento que falta, que esta en curso y que ya se completo.\\
\textit{Dado} que existen tareas en distintos estados, \textit{cuando} filtro la lista por el estado pendiente, \textit{entonces} el sistema solo muestra las tareas que estan en ese estado.}

\hu{HU-10}{RF-19}{Como administrador o tecnico, quiero generar reportes de produccion, costos y perdidas con graficos, para presentar la informacion de la finca de forma visual y tomar decisiones mas rapido.\\
\textit{Dado} un rango de fechas y uno o varios lotes seleccionados, \textit{cuando} genero el reporte, \textit{entonces} los valores mostrados corresponden exactamente a la suma de los registros de ese rango.}

\hu{HU-11}{RF-20}{Como tecnico, quiero que el sistema pida usuario y contrasena y respete los permisos de cada rol, para que cada persona solo pueda ver y hacer lo que corresponde a su funcion.\\
\textit{Dado} un usuario con rol de jornalero, \textit{cuando} intenta acceder al modulo financiero, \textit{entonces} el sistema le niega el acceso.}

\hu{HU-12}{RF-22}{Como administrador, quiero registrar el consentimiento de cada trabajador para el tratamiento de sus datos personales, para cumplir con la Ley Organica de Proteccion de Datos Personales antes de guardar cualquier informacion suya en el sistema.\\
\textit{Dado} un trabajador que va a usar el sistema por primera vez, \textit{cuando} intenta iniciar sesion sin haber registrado su consentimiento, \textit{entonces} el sistema le solicita aceptarlo antes de continuar.}

\hu{HU-13}{RF-23}{Como trabajador, quiero poder acceder, rectificar o eliminar mis propios datos personales dentro del sistema, para ejercer los derechos que me reconoce la Ley Organica de Proteccion de Datos Personales.\\
\textit{Dado} que estoy autenticado en el sistema, \textit{cuando} solicito ver, corregir o eliminar mis datos personales, \textit{entonces} el sistema atiende la solicitud y deja un registro de la accion realizada.}

\hu{HU-14}{RF-24}{Como administrador, quiero registrar el certificado de salud vigente del trabajador que aplica agroquimicos, para cumplir con lo exigido por la Resolucion 183 de AGROCALIDAD antes de asignarle esa labor.\\
\textit{Dado} un trabajador sin certificado de salud vigente registrado, \textit{cuando} intento asignarle una tarea de aplicacion de agroquimicos, \textit{entonces} el sistema advierte que el certificado no esta registrado.}

\hu{HU-15}{RF-26}{Como tecnico o jornalera, quiero registrar manualmente cada ingreso y egreso con su concepto, monto y fecha, para llevar un control independiente ademas del calculo automatico de ganancia.\\
\textit{Dado} que registro un ingreso o un egreso con su concepto, monto y fecha, \textit{cuando} reviso el reporte financiero del periodo correspondiente, \textit{entonces} ese movimiento aparece reflejado en el.}

\hu{HU-16}{RF-34}{Como tecnico, quiero que el catalogo de cultivos muestre un conjunto de campos distinto segun la variedad seleccionada, para registrar los datos agronomicos propios de cada variedad de cacao o de platano.\\
\textit{Dado} que selecciono una variedad de cacao en el formulario, \textit{cuando} el formulario se despliega, \textit{entonces} muestra un conjunto de campos distinto al que se muestra al seleccionar platano.}

\hu{HU-17}{RF-37}{Como jornalero, tecnico o administrador, quiero poder registrar y avisar de inmediato ante sintomas sospechosos de una plaga cuarentenaria como el Moko, para activar el protocolo de bioseguridad exigido por AGROCALIDAD sin perder tiempo.\\
\textit{Dado} que un usuario reporta sintomas compatibles con una plaga cuarentenaria en un lote, \textit{cuando} registra el aviso, \textit{entonces} el sistema lo marca con prioridad alta y lo deja visible para el administrador de forma inmediata.}

\subsection{Trazabilidad legal}

La cobertura de los 26 criterios legales frente a los requisitos elicitados por entrevista y frente a los requisitos derivados directamente del texto legal se documenta fila por fila en \ruta{04\_Trazabilidad/Matriz\_Trazabilidad\_v2.xlsx}. Esa matriz es la que alimenta el diseno de la prueba de McNemar del protocolo experimental, al marcar cada uno de los 26 criterios como cubierto o no cubierto antes y despues de aplicar el metodo legal-first.

\section{Modelado del sistema}

El modelado completo esta en \ruta{03\_Modelado/}, con 32 diagramas en formato \texttt{.drawio} y su exportacion a imagen, mas la especificacion textual de los casos de uso y las historias de usuario. Toma como base el diagrama de contexto de la Seccion 2.2 y el modelado organizacional i estrella de la Seccion 2.4, y desarrolla a partir de los 39 requisitos funcionales y 21 no funcionales de la Seccion 3.

\subsection{Actores y casos de uso}

El sistema tiene tres actores humanos, administrador, tecnico y trabajador agricola, y dos actores externos, el motor de recomendaciones de inteligencia artificial y AGROCALIDAD como entidad reguladora. El diagrama general de casos de uso reune 14 casos de uso en alcance, CU-01 a CU-14, desarrollados de forma completa a continuacion con precondiciones, poscondiciones, flujo principal, flujos alternativos y reglas de negocio; la fuente editable de los diagramas esta en \ruta{03\_Modelado/00\_Use\_Case\_Specifications.md} y en los 32 diagramas \texttt{.drawio} de \ruta{03\_Modelado/}. CU-15, la via de integracion con sensores del RF-32, queda documentado como Won't have y no se desarrolla.

\begin{longtable}{L{6cm} L{9cm}}
\toprule
\textbf{Caso de uso} & \textbf{Requisitos que agrupa} \\
\midrule
\endhead
CU-01 Gestionar parcelas y lotes & RF-01, RF-02, RF-16, RF-29, RF-34, RF-35 \\
CU-02 Registrar actividades agricolas & RF-03 \\
CU-03 Registrar cosecha & RF-04, RF-13, RF-14, RF-15 \\
CU-04 Calcular rendimiento y finanzas & RF-05, RF-06, RF-26 \\
CU-05 Gestionar inventario de insumos & RF-07, RF-08, RF-30 \\
CU-06 Alertas de plagas asistidas por IA & RF-09, RF-10, RF-31, RF-33 \\
CU-07 Asignar y dar seguimiento a tareas & RF-11, RF-12, RF-27, RF-28 \\
CU-08 Generar reportes & RF-19 \\
CU-09 Autenticacion y control de acceso & RF-20, RF-22, RF-23 \\
CU-10 Visitas tecnicas & RF-21 \\
CU-11 Riesgo laboral y equipo de proteccion & RF-18, RF-24 \\
CU-12 Canal rapido de reporte & RF-17 \\
CU-13 Cumplimiento BPA & RF-25, RF-36, RF-38, RF-39 \\
CU-14 Aviso de plaga cuarentenaria & RF-37, RF-38 \\
\bottomrule
\end{longtable}

\subsubsection{CU-01. Gestionar parcelas y lotes}
\begin{itemize}
\item \textbf{Actor(es) principal(es):} Administrador, Tecnico, Jornalero.
\item \textbf{Requisitos relacionados:} RF-01, RF-02, RF-16, RF-29, RF-34, RF-35.
\item \textbf{Precondicion:} usuario autenticado con permiso de edicion (RF-01); catalogo de cultivos ya configurado (RF-02).
\item \textbf{Flujo principal:}
  \begin{enumerate}
  \item El usuario selecciona "Nueva parcela".
  \item Ingresa nombre, sector, ubicacion, area, cultivo, variedad y cantidad de plantas.
  \item El sistema muestra el conjunto de campos correspondiente a la variedad seleccionada (RF-34).
  \item El usuario guarda la parcela.
  \item El sistema muestra la parcela en el listado general en menos de 2 segundos.
  \end{enumerate}
\item \textbf{Flujos alternativos:}
  \begin{itemize}
  \item 2a. El usuario deja pendiente la condicion climatica o de suelo: puede registrarla en un momento posterior, sin que eso bloquee el guardado de la parcela (RF-16).
  \item 4a. El usuario registra un analisis de suelo previo a la siembra desde la misma parcela, en cualquier momento (RF-35).
  \end{itemize}
\item \textbf{Excepciones:}
  \begin{itemize}
  \item 2e. El usuario intenta ingresar un cultivo fuera del catalogo cerrado: el sistema lo rechaza e indica que debe elegir uno de los cultivos configurados (RF-02).
  \item 4e. Falta alguno de los 7 campos obligatorios: el sistema rechaza el guardado y marca los campos vacios (RF-01).
  \item 5e. La parcela no tiene ningun evento registrado al consultar su historial: el sistema muestra el historial vacio, no un error (RF-29).
  \end{itemize}
\item \textbf{Postcondicion:} parcela visible en el listado general (RF-01); ningun registro de cultivo queda con valor "otros" sin especificar (RF-02); historial de la parcela consultable desde su creacion (RF-29).
\item \textbf{Reglas de negocio:} CU-01 no permite valor "otros" para cultivo (RF-02); el catalogo de cultivo/variedad solo lo edita el Administrador; una parcela marcada "para exportacion" habilita los campos de ExportBatch (RF-14, ver CU-03).
\end{itemize}

\subsubsection{CU-02. Registrar actividades agricolas}
\begin{itemize}
\item \textbf{Actor(es) principal(es):} Jornalero, Tecnico.
\item \textbf{Requisitos relacionados:} RF-03, RNF-04.
\item \textbf{Precondicion:} la parcela ya existe (RF-01 ya ejecutado).
\item \textbf{Flujo principal:}
  \begin{enumerate}
  \item El usuario selecciona la parcela.
  \item Ingresa fecha, tipo de actividad (fumigacion, fertilizacion, poda, riego, limpieza), insumo/producto usado y trabajador responsable.
  \item El sistema guarda el registro y lo vincula al historial de la parcela.
  \end{enumerate}
\item \textbf{Flujos alternativos:}
  \begin{itemize}
  \item 3a. Sin conexion al momento de guardar: el sistema almacena el registro en el dispositivo y sincroniza al reconectar, sin perdida de datos dentro del limite de 7 dias sin conexion (RNF-04).
  \end{itemize}
\item \textbf{Excepciones:}
  \begin{itemize}
  \item 2e. Falta fecha, parcela o trabajador responsable: el sistema rechaza el guardado (RF-03).
  \end{itemize}
\item \textbf{Postcondicion:} la actividad queda vinculada al historial de la parcela.
\item \textbf{Reglas de negocio:} toda actividad que use un insumo del catalogo decrementa automaticamente el inventario (ver CU-05); el historial de actividades es la fuente del "historial de tratamientos" que consume CU-06 (Alertas IA).
\end{itemize}

\subsubsection{CU-03. Registrar cosecha}
\begin{itemize}
\item \textbf{Actor(es) principal(es):} Jornalero.
\item \textbf{Requisitos relacionados:} RF-04, RF-13, RF-14, RF-15, RNF-04.
\item \textbf{Precondicion:} parcela y cultivo existen (RF-04); cosecha ya registrada ese dia (para RF-13).
\item \textbf{Flujo principal:}
  \begin{enumerate}
  \item El usuario selecciona parcela y cultivo.
  \item El sistema determina la unidad esperada para el cultivo (baba/quintal para cacao, racimo/caja para platano).
  \item El usuario ingresa cantidad y unidad.
  \item El sistema valida la unidad, guarda la cosecha y actualiza el total de parcela/periodo.
  \item Si aplica, el usuario consulta el precio de mercado actual y el sistema estima el ingreso proyectado (RF-15).
  \end{enumerate}
\item \textbf{Flujos alternativos:}
  \begin{itemize}
  \item 3a. El usuario registra una cantidad rechazada y su motivo (enfermedad, plaga, dano mecanico): el sistema excluye automaticamente esa cantidad del total utilizable (RF-13).
  \item 3b. La parcela es de exportacion: el usuario registra color de cinta, semana de enfunde y calidad de caja (RF-14).
  \item 4a. Sin conexion al guardar: el sistema almacena la cosecha localmente y sincroniza al reconectar, dentro del limite de 7 dias sin perdida de datos (RNF-04).
  \item 5a. No hay un precio de mercado disponible de fuente externa: el usuario lo ingresa manualmente para obtener la estimacion (RF-15).
  \end{itemize}
\item \textbf{Excepciones:}
  \begin{itemize}
  \item 3e. No hay una cosecha registrada ese dia: el sistema no permite registrar un rechazo asociado (RF-13).
  \item 4e. Unidad de cosecha invalida para el cultivo seleccionado: el sistema rechaza el registro (RF-04).
  \end{itemize}
\item \textbf{Postcondicion:} cosecha sumada al total de la parcela/periodo; historial completo de empaque a exportacion consultable si la parcela es de exportacion (RF-14).
\item \textbf{Reglas de negocio:} la unidad de cosecha depende exclusivamente del cultivo, nunca de texto libre; el total utilizable siempre excluye el rechazo registrado.
\end{itemize}

\subsubsection{CU-04. Calcular rendimiento y finanzas}
\begin{itemize}
\item \textbf{Actor(es) principal(es):} Administrador, Tecnico, Jornalero.
\item \textbf{Requisitos relacionados:} RF-05, RF-06, RF-26.
\item \textbf{Precondicion:} existen registros de ingresos y egresos del periodo (RF-06).
\item \textbf{Flujo principal:}
  \begin{enumerate}
  \item El sistema consulta el historial de cosecha de la parcela y calcula el rendimiento promedio por periodo, permitiendo comparar periodos.
  \item El sistema consulta los registros de ingresos y egresos (automaticos por insumos y manuales por RF-26) del periodo.
  \item El sistema calcula la ganancia neta (ingresos menos egresos) y la muestra en el modulo financiero.
  \end{enumerate}
\item \textbf{Flujos alternativos:}
  \begin{itemize}
  \item 2a. El usuario registra manualmente un ingreso o egreso adicional (concepto, monto, fecha, tipo), independiente del calculo automatico (RF-26).
  \end{itemize}
\item \textbf{Excepciones:}
  \begin{itemize}
  \item 1e. Existen menos de 2 registros de cosecha para la parcela: el sistema muestra "datos insuficientes" en vez de un promedio (RF-05).
  \item 2e. No existen registros de ingreso o egreso en el periodo consultado: el sistema muestra el modulo financiero en cero, sin marcar error (RF-06).
  \end{itemize}
\item \textbf{Postcondicion:} valor mostrado en el reporte de parcela (RF-05); valor visible en el modulo financiero (RF-06); movimiento reflejado en el calculo de ganancia neta (RF-26).
\item \textbf{Reglas de negocio:} ganancia neta = suma de ingresos menos suma de egresos del periodo, sin excepcion; el rendimiento promedio se calcula como la suma dividida entre N registros.
\end{itemize}

\subsubsection{CU-05. Gestionar inventario de insumos}
\begin{itemize}
\item \textbf{Actor(es) principal(es):} Administrador, Tecnico, Jornalero.
\item \textbf{Requisitos relacionados:} RF-07, RF-08, RF-30, RNF-15.
\item \textbf{Precondicion:} un umbral definido para el insumo (RF-08).
\item \textbf{Flujo principal:}
  \begin{enumerate}
  \item El usuario registra el uso de una cantidad de un insumo.
  \item El sistema actualiza el stock disponible.
  \item El sistema verifica si el nuevo stock esta por debajo del umbral minimo configurado (RF-30).
  \item Si es asi, el sistema notifica al responsable en menos de 5 minutos.
  \end{enumerate}
\item \textbf{Flujos alternativos:}
  \begin{itemize}
  \item 3a. El Tecnico configura o edita el umbral minimo de stock de un insumo en cualquier momento; el nuevo valor se aplica desde la siguiente actualizacion de stock (RF-30).
  \item 4a. Conectividad intermitente al momento de notificar: el sistema reintenta la entrega y la confirma dentro de los 5 minutos posteriores a la reconexion del dispositivo (RNF-15).
  \end{itemize}
\item \textbf{Excepciones:}
  \begin{itemize}
  \item 2e. El uso registrado dejaria el stock en un valor negativo: el sistema advierte antes de permitir guardar.
  \item 3e. El insumo todavia no tiene un umbral configurado: el sistema no dispara ninguna alerta hasta que el Tecnico lo defina (RF-08).
  \end{itemize}
\item \textbf{Postcondicion:} stock reducido tras cada uso registrado (RF-07); notificacion enviada y registrada si se cruza el umbral (RF-08).
\item \textbf{Reglas de negocio:} CU-05 dispara una alerta al cruzar el stock minimo (RF-08, RF-30); el stock nunca puede quedar negativo.
\end{itemize}

\subsubsection{CU-06. Alertas de plagas asistidas por IA}
\begin{itemize}
\item \textbf{Actor(es) principal(es):} Administrador, Tecnico, Jornalero.
\item \textbf{Actor(es) secundario(s):} Sistema de IA.
\item \textbf{Requisitos relacionados:} RF-09, RF-10, RF-31, RF-33, RNF-16, RNF-17, RNF-18, RNF-19, RNF-20, RNF-21.
\item \textbf{Precondicion:} modelo de IA entrenado disponible (RF-09); conexion a internet (RF-10, si se usa diagnostico por imagen).
\item \textbf{Flujo principal:}
  \begin{enumerate}
  \item El Sistema de IA analiza los datos de la parcela (historial de tratamientos y, opcionalmente, una foto subida por el usuario) y genera una sugerencia con justificacion y fuente de datos (RNF-17).
  \item El sistema notifica al usuario responsable, marcada explicitamente como sugerencia por confirmar en campo.
  \item El usuario revisa la justificacion y confirma o descarta la sugerencia.
  \item El sistema registra la decision.
  \end{enumerate}
\item \textbf{Flujos alternativos:}
  \begin{itemize}
  \item 1a. El sistema compara un valor recien ingresado contra el promedio historico de la parcela (minimo 3 registros previos) y advierte si se desvia mas de 2 desviaciones estandar, antes de guardar (RF-31).
  \item 3a. El usuario descarta la sugerencia de IA (RF-09).
  \end{itemize}
\item \textbf{Excepciones:}
  \begin{itemize}
  \item 1e. No hay conexion a internet al intentar el diagnostico por imagen: el sistema informa que esa funcion no esta disponible en ese momento (RF-10).
  \item 1f. La parcela no tiene historico suficiente (menos de 3 registros): el sistema guarda el dato sin comparar contra un promedio (RF-31).
  \item 3e. El usuario esta en desacuerdo con la recomendacion mas alla de un simple descarte: registra formalmente su desacuerdo, disponible para revision del equipo tecnico, y la alerta queda fuera del conjunto aplicado hasta esa revision (RF-33).
  \end{itemize}
\item \textbf{Postcondicion:} la alerta queda registrada con el resultado de la verificacion humana (RF-09).
\item \textbf{Reglas de negocio:} ninguna alerta se aplica automaticamente sobre una parcela sin confirmacion del usuario responsable (RF-09, RNF-19); toda alerta muestra su justificacion (maximo 60 palabras) y al menos una fuente de datos antes del boton de confirmacion (RNF-17).
\end{itemize}

\subsubsection{CU-07. Asignar y dar seguimiento a tareas}
\begin{itemize}
\item \textbf{Actor(es) principal(es):} Administrador, Tecnico, Jornalero.
\item \textbf{Requisitos relacionados:} RF-11, RF-12, RF-27, RF-28, RNF-15.
\item \textbf{Precondicion:} trabajador registrado (RF-11); trabajador con dispositivo registrado (RF-12).
\item \textbf{Flujo principal:}
  \begin{enumerate}
  \item El usuario asigna una tarea a un trabajador, indicando parcela, fecha y tipo.
  \item El sistema crea la tarea con estado "pendiente" y la muestra con su estado.
  \item El sistema notifica al trabajador (push/SMS) incluyendo parcela, fecha y tipo de tarea.
  \item El trabajador actualiza el estado de la tarea (pendiente, en progreso, bloqueada, completada) y puede adjuntar una observacion en texto libre (RF-28).
  \item El trabajador puede consultar en cualquier momento la vista de sus propias tareas pendientes (RF-27).
  \end{enumerate}
\item \textbf{Flujos alternativos:}
  \begin{itemize}
  \item 1a. La parcela tiene un riesgo laboral registrado (ver CU-11): el sistema muestra la advertencia de EPP antes de confirmar la asignacion.
  \item 3a. Conectividad intermitente al notificar: el sistema reintenta la entrega y confirma dentro de los 5 minutos posteriores a la reconexion (RNF-15).
  \end{itemize}
\item \textbf{Excepciones:}
  \begin{itemize}
  \item 3e. El trabajador no tiene un dispositivo registrado: el sistema deja la tarea visible en su vista de pendientes (RF-27) aunque no pueda enviar la notificacion push (RF-12).
  \end{itemize}
\item \textbf{Postcondicion:} tarea filtrable por sus 4 estados (RF-11); notificacion entregada (RF-12).
\item \textbf{Reglas de negocio:} el filtro de estado muestra exclusivamente las tareas en ese estado, sin mezclar parcelas o trabajadores; la observacion de una tarea completada permanece visible al ver el detalle de la tarea (RF-28).
\end{itemize}

\subsubsection{CU-08. Generar reportes}
\begin{itemize}
\item \textbf{Actor(es) principal(es):} Administrador, Tecnico.
\item \textbf{Requisitos relacionados:} RF-19.
\item \textbf{Precondicion:} existen datos en el rango seleccionado.
\item \textbf{Flujo principal:}
  \begin{enumerate}
  \item El usuario selecciona un rango de fechas y uno o mas lotes.
  \item El sistema consolida produccion, costos, perdidas por plaga y rendimiento del rango seleccionado.
  \item El sistema genera graficos de barra y circulares codificados por color, reflejando exactamente la suma de los registros del rango.
  \end{enumerate}
\item \textbf{Flujos alternativos:} ninguno adicional a los cubiertos por RF-19.
\item \textbf{Excepciones:}
  \begin{itemize}
  \item 2e. No existen datos en el rango de fechas seleccionado: el sistema muestra "no hay datos disponibles" en vez de un reporte vacio (RF-19).
  \end{itemize}
\item \textbf{Postcondicion:} reporte generado y descargable.
\item \textbf{Reglas de negocio:} el reporte generado refleja exactamente la suma de los registros del rango seleccionado, sin redondeo no declarado.
\end{itemize}

\subsubsection{CU-09. Autenticacion y control de acceso}
\begin{itemize}
\item \textbf{Actor(es) principal(es):} Administrador, Tecnico, Jornalero.
\item \textbf{Requisitos relacionados:} RF-20, RF-22, RF-23.
\item \textbf{Precondicion:} usuario previamente registrado (RF-20).
\item \textbf{Flujo principal:}
  \begin{enumerate}
  \item El usuario ingresa usuario y contrasena.
  \item El sistema valida las credenciales y determina el rol.
  \item Si es el primer inicio de sesion, el sistema muestra el aviso de tratamiento de datos personales (LOPDP Art. 8) y solicita la aceptacion explicita del usuario antes de continuar.
  \item El sistema registra el consentimiento (libre, especifico, informado e inequivoco) y otorga acceso segun el rol.
  \end{enumerate}
\item \textbf{Flujos alternativos:}
  \begin{itemize}
  \item 4a. El trabajador solicita acceder, rectificar o eliminar sus propios datos personales en cualquier momento posterior, desde Ajustes -> Mis Datos (derechos ARCO+, RF-23).
  \end{itemize}
\item \textbf{Excepciones:}
  \begin{itemize}
  \item 2e. Credenciales invalidas: el sistema niega el acceso.
  \item 3e. El usuario no acepta el aviso de tratamiento de datos: el sistema no concede acceso a ningun modulo que trate datos personales (RF-22).
  \item 4e. Un usuario con rol "jornalero" intenta acceder al modulo financiero: el sistema niega el acceso (RF-20).
  \item 4f. La solicitud ARCO+ es de eliminacion y existen datos que deben conservarse por obligacion legal (por ejemplo, trazabilidad BPA): el sistema informa cuales datos no pueden eliminarse y por que (RF-23).
  \end{itemize}
\item \textbf{Postcondicion:} acceso otorgado solo a las funciones del rol del usuario (RF-20); consentimiento registrado antes de cualquier tratamiento de datos personales (RF-22).
\item \textbf{Reglas de negocio:} CU-09 es un Must-have regulatorio: ningun modulo es accesible sin autenticacion y sin consentimiento LOPDP registrado; el rol determina de forma exhaustiva el conjunto de funciones visibles.
\end{itemize}

\subsubsection{CU-10. Visitas tecnicas}
\begin{itemize}
\item \textbf{Actor principal:} Tecnico.
\item \textbf{Requisitos relacionados:} RF-21.
\item \textbf{Precondicion:} ninguna previa al registro de la visita.
\item \textbf{Flujo principal:}
  \begin{enumerate}
  \item El Tecnico registra una visita tecnica: fecha, hallazgos e informe adjunto.
  \item El sistema guarda la visita y la asocia al historial de cumplimiento del lote.
  \item Transcurridos 15 dias desde la visita, el sistema genera automaticamente un recordatorio de la siguiente visita.
  \end{enumerate}
\item \textbf{Flujos alternativos:}
  \begin{itemize}
  \item 2a. Se encuentra un incumplimiento: el sistema sugiere registrar la capacitacion o el EPP relacionado (ver CU-11, CU-13).
  \item 3a. El recordatorio se genera sin intervencion del Tecnico, exactamente a los 15 dias de la visita registrada (RF-21).
  \end{itemize}
\item \textbf{Excepciones:} ninguna adicional a las cubiertas por RF-21.
\item \textbf{Postcondicion:} recordatorio generado para la siguiente visita (RF-21).
\item \textbf{Reglas de negocio:} las visitas tecnicas periodicas son la fuente principal de evidencia para la renovacion de la certificacion BPA (CU-13).
\end{itemize}

\subsubsection{CU-11. Riesgo laboral y equipo de proteccion}
\begin{itemize}
\item \textbf{Actor(es) principal(es):} Jornalero, Tecnico.
\item \textbf{Requisitos relacionados:} RF-18, RF-24.
\item \textbf{Precondicion:} la parcela existe; hay una tarea por asignar o en curso.
\item \textbf{Flujo principal:}
  \begin{enumerate}
  \item Al asignar o recibir una tarea, el sistema verifica si la parcela tiene un riesgo laboral registrado.
  \item Si no esta registrado, el Tecnico registra tipo de riesgo, severidad y EPP sugerido.
  \item El sistema muestra la advertencia de EPP al trabajador antes de confirmar la asignacion.
  \end{enumerate}
\item \textbf{Flujos alternativos:}
  \begin{itemize}
  \item 2a. El Administrador adjunta el certificado de salud del trabajador al registro (RF-24, Res. AGROCALIDAD 183 Art. 33-34).
  \end{itemize}
\item \textbf{Excepciones:}
  \begin{itemize}
  \item 2e. La tarea es de aplicacion de agroquimicos y el trabajador no tiene certificado de salud vigente registrado: el sistema advierte que el certificado no esta registrado al intentar asignarle esa tarea (RF-24).
  \item 3e. La tarea se intenta confirmar sin que se haya mostrado la advertencia de EPP: el sistema bloquea la confirmacion.
  \end{itemize}
\item \textbf{Postcondicion:} riesgo registrado con EPP sugerido (RF-18); certificado de salud adjunto cuando aplica (RF-24).
\item \textbf{Reglas de negocio:} ninguna tarea sobre una parcela con riesgo registrado puede confirmarse sin que se muestre primero la advertencia de EPP.
\end{itemize}

\subsubsection{CU-12. Canal rapido de reporte}
\begin{itemize}
\item \textbf{Actor principal:} Jornalero.
\item \textbf{Requisitos relacionados:} RF-17, RNF-03.
\item \textbf{Precondicion:} usuario autenticado.
\item \textbf{Flujo principal:}
  \begin{enumerate}
  \item El trabajador abre el canal de reporte rapido.
  \item Graba un mensaje breve (texto o voz) describiendo una situacion de campo.
  \item El sistema almacena el reporte y lo marca para revision/triage del Tecnico, en un maximo de 2 toques desde la pantalla principal.
  \end{enumerate}
\item \textbf{Flujos alternativos:}
  \begin{itemize}
  \item 2a. El reporte describe una posible plaga cuarentenaria: el sistema sugiere escalarlo como aviso fitosanitario formal (ver CU-14, RF-37).
  \end{itemize}
\item \textbf{Excepciones:} ninguna adicional a las cubiertas por RF-17.
\item \textbf{Postcondicion:} reporte rapido (chat/voz) registrado y enrutado al modulo correspondiente (actividad, enfermedad, riesgo).
\item \textbf{Reglas de negocio:} el canal rapido de reporte favorece la baja alfabetizacion digital (RNF-03) y no requiere campos estructurados para iniciar.
\end{itemize}

\subsubsection{CU-13. Cumplimiento BPA}
\begin{itemize}
\item \textbf{Actor principal:} Administrador.
\item \textbf{Requisitos relacionados:} RF-25, RF-36, RF-38, RF-39.
\item \textbf{Precondicion:} visitas tecnicas y registros de capacitacion disponibles como evidencia de soporte.
\item \textbf{Flujo principal:}
  \begin{enumerate}
  \item El Administrador revisa el tablero de cumplimiento (vigencia del certificado, capacitaciones pendientes, completitud de la bitacora de bioseguridad).
  \item Si el certificado esta por vencer o falta evidencia de soporte, el Administrador solicita la renovacion ante AGROCALIDAD.
  \item El sistema adjunta los registros de capacitacion (manejo de plaguicidas, primeros auxilios) y la bitacora de ingreso/salida de bioseguridad como evidencia de soporte.
  \end{enumerate}
\item \textbf{Flujos alternativos:}
  \begin{itemize}
  \item 1a. Cada visita confirmada en CU-10 o CU-14 genera automaticamente un registro en la bitacora de bioseguridad de ingreso/salida, sin que el Administrador deba capturarlo por separado (RF-38).
  \end{itemize}
\item \textbf{Excepciones:}
  \begin{itemize}
  \item 2e. Falta una capacitacion obligatoria: el sistema bloquea la solicitud de renovacion hasta que se registre (RF-36, RF-39).
  \end{itemize}
\item \textbf{Postcondicion:} solicitud/renovacion de certificacion BPA presentada ante AGROCALIDAD (RF-25); registros de capacitacion y bioseguridad disponibles para auditoria (RF-36, RF-38, RF-39).
\item \textbf{Reglas de negocio:} este CU es un Must-have regulatorio (Res. AGROCALIDAD 183); su ausencia era un vacio legal identificado por el metodo legal-first del equipo.
\end{itemize}

\subsubsection{CU-14. Aviso de plaga cuarentenaria (p. ej. Moko)}
\begin{itemize}
\item \textbf{Actor(es) principal(es):} Jornalero, Tecnico, Administrador.
\item \textbf{Actor(es) secundario(s):} AGROCALIDAD.
\item \textbf{Requisitos relacionados:} RF-37, RF-38.
\item \textbf{Precondicion:} sintomas sospechosos de plaga cuarentenaria detectados en una parcela (Res. AGROCALIDAD 0072, Art. 3.6.1.a).
\item \textbf{Flujo principal:}
  \begin{enumerate}
  \item Un usuario de campo detecta sintomas sospechosos (p. ej. Moko) en una parcela.
  \item Registra el sintoma y la parcela afectada en el sistema.
  \item El sistema genera un aviso fitosanitario con prioridad alta, visible de inmediato para el Administrador.
  \item Un Tecnico confirma el sintoma tras la inspeccion.
  \item El sistema deja el aviso disponible para presentacion formal ante AGROCALIDAD.
  \end{enumerate}
\item \textbf{Flujos alternativos:}
  \begin{itemize}
  \item 4a. Al confirmarse el sintoma, el sistema registra automaticamente el evento en la bitacora de bioseguridad de ingreso/salida de la visita, sin que el usuario deba capturarlo por separado (RF-38).
  \end{itemize}
\item \textbf{Excepciones:}
  \begin{itemize}
  \item 4e. El Tecnico descarta el sintoma tras la inspeccion: el aviso no se envia a AGROCALIDAD, pero queda registrado como descartado, para trazabilidad.
  \end{itemize}
\item \textbf{Postcondicion:} aviso registrado y disponible para presentacion formal ante AGROCALIDAD.
\item \textbf{Reglas de negocio:} este CU es un Must-have regulatorio (Res. AGROCALIDAD 0072, Art. 3.6.1.a); su ausencia era un vacio legal identificado por el metodo legal-first del equipo; ningun aviso se envia a AGROCALIDAD sin la confirmacion de un Tecnico.
\end{itemize}

Las historias de usuario con criterios de aceptacion en Gherkin estan en \ruta{03\_Modelado/00\_User\_Stories\_Acceptance\_Criteria.md}, una por cada requisito funcional Must have que la lleva, segun la regla de la guia de modelado.

\subsection{Modelo de dominio}

El dominio se reparte en dos diagramas de clases. \texttt{CD01\_Refined\_Class\_Diagram} cubre el dominio operativo, con la jerarquia de usuario, parcela, cultivo, cosecha, labor, insumo, reporte y las clases de recomendacion y alerta de inteligencia artificial. \texttt{CD02\_Legal\_Compliance\_Class\_Diagram} cubre el dominio de cumplimiento legal derivado del enfoque legal-first, con las clases de consentimiento, solicitud de derechos del titular, visita tecnica, riesgo laboral, aviso fitosanitario, bitacora de bioseguridad, certificacion BPA, registro de capacitacion, lote de exportacion y transaccion financiera.

\subsection{Comportamiento}

El comportamiento se documenta con 7 diagramas de actividad, con calles por actor, que cubren el flujo operativo principal, la alerta de inventario, las recomendaciones de inteligencia artificial, el consentimiento y los derechos ARCO+ de la LOPDP, la visita tecnica y el cumplimiento BPA, el riesgo laboral y el aviso de plaga cuarentenaria. Los 14 diagramas de secuencia siguen la separacion en capas de frontera, control y entidad, con fragmentos de alternativa y barras de activacion. Dos diagramas de estados modelan el ciclo de vida de una tarea de labor y el de una alerta de inteligencia artificial. Este segundo diagrama fija la regla central del RF-09, ninguna alerta se aplica de forma automatica a un lote sin la confirmacion del usuario responsable, y el camino del RF-33, en el que un desacuerdo formal mantiene la alerta fuera del conjunto aplicado hasta que la revise el equipo tecnico.

\subsection{Arquitectura de componentes}

El diagrama de componentes divide el backend en modulos, \texttt{AuthModule} y \texttt{UserModule} para autenticacion y usuarios, \texttt{PlotCropModule} para parcelas y catalogo, \texttt{TaskModule} para tareas, \texttt{InventoryModule} para insumos, \texttt{ReportModule} para cosecha, rendimiento y reportes, \texttt{AIIntegrationModule} para la integracion con el motor de recomendaciones, y \texttt{ComplianceModule} para todo el dominio de cumplimiento LOPDP y AGROCALIDAD. El cliente es una aplicacion web que consume una API REST sobre HTTPS. Los actores externos son el motor de recomendaciones, conectado por su propia API sin capacidad de modificar datos de forma autonoma, y la pasarela de reporte a AGROCALIDAD para los avisos de las Resoluciones 183 y 0072.

\subsection{Mockups}

Hay 9 pantallas en \ruta{03\_Modelado/Mockups/}, MU-01 a MU-09, cada una en un archivo HTML autonomo y exportada a imagen, mas \texttt{MU-00\_Prototype.html} como prototipo navegable que las enlaza. Cubren inicio de sesion y consentimiento, listado de parcelas, registro de actividad, registro de cosecha, alerta de inventario, alerta de inteligencia artificial, tablero de tareas, panel de reportes y la pantalla de cumplimiento y visitas tecnicas.

\subsection{Diagrama de despliegue}

El conjunto se levanta con \texttt{docker compose up} desde \ruta{05\_MVP/}. El backend expone la logica de negocio a traves de una API REST y la base de datos PostgreSQL cuenta con un chequeo de salud antes de aceptar trafico del backend.

\begin{verbatim}
        +---------------------------------------------+
        |               Host Docker                   |
        |  +---------------------+  +---------------+  |
        |  | agromoreira_backend |->| agromoreira_db|  |
        |  | Node.js / Express   |  | PostgreSQL 16 |  |
        |  | puerto 3000         |  | puerto 5432   |  |
        |  +---------------------+  +---------------+  |
        +----------^-------------------------|--------+
                   |                         |
             cliente web o movil     volumen persistente
                (HTTP)               agromoreira_db_data
\end{verbatim}

El diagrama de despliegue completo, con el dispositivo cliente, el servidor de aplicacion, el servidor de base de datos y los nodos externos del motor de recomendaciones y de AGROCALIDAD, esta en \ruta{03\_Modelado/Diagrams/DEP01\_Deployment\_Diagram.drawio}.

\section{Priorizacion y trazabilidad}

Los 39 requisitos funcionales cuentan con su prioridad MoSCoW asignada dentro de cada ficha de la Seccion 3. Sobre esa base, el equipo completo realizo el 1 de septiembre de 2026 una sesion de priorizacion complementaria con el modelo de Kano y el metodo WSJF, del marco agil SAFe. La clasificacion Kano de cada requisito, en basico, de desempeno, atractivo o indiferente, se definio por consenso del equipo a partir de la evidencia cualitativa recogida en las 17 entrevistas y no de una encuesta Kano formal aplicada a los participantes, ya que esa encuesta especifica no formo parte del instrumento de campo. Para el WSJF, el equipo asigno a cada requisito cuatro valores del 1 al 10, valor de negocio, urgencia, reduccion de riesgo u oportunidad, y tamano del esfuerzo, y el puntaje final se calculo como la suma de los tres primeros valores dividida entre el tamano del esfuerzo.

El resultado completo, con los 39 requisitos y su ranking final, se documenta en \ruta{priorizacion\_moscow\_kano.csv}. Los diez requisitos de mayor prioridad segun este metodo son los siguientes.

\begin{longtable}{L{1.4cm} L{7.2cm} L{3.2cm} L{2cm} L{1.2cm}}
\toprule
\textbf{Ranking} & \textbf{Requisito} & \textbf{MoSCoW} & \textbf{Kano} & \textbf{WSJF} \\
\midrule
\endhead
1 & RF-26, registro manual de ingresos y egresos & Must have & Basico & 6.5 \\
2 & RF-08, alertas automaticas de bajo stock & Must have & De desempeno & 6.3 \\
2 & RF-13, registro de motivo de rechazo de fruta & Should have & De desempeno & 6.3 \\
4 & RF-02, catalogo cerrado de cultivos y variedades & Should have & Basico & 6.0 \\
4 & RF-22, consentimiento del trabajador para el tratamiento de datos & Must have regulatorio & Basico & 6.0 \\
4 & RF-37, aviso ante plaga cuarentenaria & Must have regulatorio & Basico & 6.0 \\
7 & RF-18, registro de riesgo laboral y equipo de proteccion & Should have & Basico & 5.6 \\
7 & RF-24, certificado de salud del trabajador & Must have regulatorio & Basico & 5.6 \\
7 & RF-38, bitacora de bioseguridad de ingreso y salida & Should have regulatorio & Basico & 5.6 \\
10 & RF-30, umbral de stock minimo configurable & Should have & De desempeno & 5.3 \\
\bottomrule
\end{longtable}

El requisito con menor prioridad resultante es el RF-32, la via de integracion futura con sensores de campo, coherente con su prioridad MoSCoW de Won't have para esta version.

La matriz de trazabilidad esta cerrada en \ruta{04\_Trazabilidad/Matriz\_Trazabilidad\_v2.xlsx}, con 60 filas. Son 54 filas base, una por cada requisito funcional, no funcional o derivado del catalogo actual, y 6 filas secundarias para los seis requisitos que dan servicio a un segundo caso de uso ademas del principal. Cada fila enlaza la ley o la fuente de elicitacion, el articulo, el objetivo, el interesado, el codigo de evidencia, el requisito, el tipo, el caso de uso, la historia de usuario y su criterio de aceptacion cuando corresponde, el componente y el mockup. Todos los identificadores de caso de uso, componente y mockup coinciden con los del paquete de modelado de la Seccion 4.

El tratamiento de huerfanos y cadenas rotas se documenta en \ruta{04\_Trazabilidad/Analisis\_Huerfanos\_Cadenas\_Rotas.md} y se verifica de forma automatizada con \ruta{04\_Trazabilidad/verificar\_trazabilidad.py}, que contrasta la matriz contra el ERS y contra la especificacion de casos de uso. La ejecucion del 5 de septiembre de 2026 no encuentra cadenas rotas ni huerfanos en ninguna de las dos direcciones: los 60 requisitos y los 15 casos de uso estan trazados y ninguna fila deja un eslabon vacio. Las 16 filas cuya evidencia de elicitacion figura como ``Sin evidencia'' se clasifican alli por causa y accion: ocho son requisitos derivados de norma legal, dos son requisitos no funcionales derivados de norma, y seis son atributos de calidad del componente de IA cuyo fundamento normativo se documenta en \ruta{01\_ERS/Clasificacion\_Riesgo\_IA.md}.


\section{Producto minimo viable}

El producto minimo viable esta en \ruta{05\_MVP/} (ver \ruta{05\_MVP/README.md} y \ruta{05\_MVP/docker-compose.yml}), servido por \texttt{nginx:1.25-alpine} via \texttt{docker compose up} (puerto 8080$\to$80, \texttt{http://localhost:8080}). El \ruta{05\_MVP/index.html} implementa 8 Must-have (RF-01, RF-03, RF-04, RF-07, RF-08, RF-11, RF-19, RF-20) y, junto con los mockups clicables de \ruta{03\_Modelado/Mockups/} (MU-01, MU-06, MU-08, MU-09), alcanza 14 de los 17 Must-have (82\%, $\geq$80\% exigido en C3): RF-01, RF-03, RF-04, RF-07, RF-08, RF-11, RF-19, RF-20 directos, mas RF-22, RF-23, RF-24, RF-26, RF-34 y RF-37 via mockups (RF-23, derechos ARCO+, se ejerce desde el aviso de consentimiento de MU-01, que lo enuncia de forma explicita; RF-02 es Should-have y solo se incluye como apoyo, no cuenta entre los 17 Must-have). Quedan fuera del prototipo actual RF-05, RF-06 y RF-09 (calculo de rendimiento/ganancia y alerta IA con verificacion humana), que requieren servicio externo y quedan como trabajo futuro. La matriz \ruta{04\_Trazabilidad/Matriz\_Trazabilidad\_v2.xlsx} traza cada RF a su CU, componente y mockup.

\bibliographystyle{plain}
\bibliography{referencias}
\addcontentsline{toc}{section}{Referencias}

\end{document}
